% !TeX root = ../../Book2.tex
% 纲要标题：## 第5章　范畴论初步
% 纲要位置：计划纲要.md，第 950 行。

\chapter{范畴论初步}
\label{b2:ch:05}
\BookChapterNumber{b2:ch:05}

\begin{chapterabstract}
本章从模的构造所满足的映射性质出发，介绍范畴、函子与自然变换。积、余积、拉回和推出在模范畴中得到具体构造和证明；范畴等价则通过全忠实与本质满等条件辨认。
\end{chapterabstract}

\begin{learninggoals}
\begin{itemize}
\item 从对象、态射和复合三个方面说明一个范畴，核对函子的协变或反变方向。
\item 写出自然性等式，区分依赖选择的单个同构、逐对象同构和自然同构。
\item 构造模的积、余积、拉回与推出，证明相容映射的存在唯一性，利用泛性质比较不同实现。
\item 掌握全、忠实、本质满的区别，完整构造范畴等价的拟逆，理解等价与严格同构的差异。
\end{itemize}
\end{learninggoals}

给一个 \(R\) 模 \(N\) 中的两个元素 \(u,v\)，便能唯一指定从自由模 \(R^2\) 到 \(N\) 的同态，使两个标准基向量分别映到 \(u,v\)。如果先在 \(R^2\) 中加入关系 \(e_1=e_2\)，同样的指定只有在 \(u=v\) 时才能定义在商模上。这种“何时可以指定映射”的问题，不必先选商模元素的代表元来表达。把它作为对象的刻画，便可比较由不同构造得到、却承担同一映射任务的对象。

\ProseParagraphBreak

相同的想法也适用于直和：分别给出两个模到目标的映射，就应当得到从直和出发的唯一映射；若目标本身是两个模的积，则要从一个对象出发分别给出映射。两个方向不能混淆。把对象和箭头一起考虑，构造的特征便能写成统一的映射条件；随后还要检查当原来的模或同态改变时，这些对应是否仍相容。

\ProseParagraphBreak

本章需要第二章的模同态、商模和直和，以及\cref{b2:thm:free-module-universal}的自由模泛性质；双对偶的例子调用\cref{b2:thm:double-dual}。矩阵与有限维向量空间之间的等价由选基和线性映射建立。抽象定义旁的模论例子显示箭头方向为何重要：拉回的映射指向公共目标，推出的映射来自公共定义域。第三章的正合列及下一章的张量积，也可以用这种箭头语言重新理解。

\ProseParagraphBreak
关于函子与自然性，可参阅 Riehl 的《Category Theory in Context》\cite[Definition 1.3.1, p. 13; Definition 1.4.1, pp. 23--24]{Riehl2016CategoryTheory}；该书的等价判据见\cite[Theorem 1.5.9, pp. 31--32]{Riehl2016CategoryTheory}。下一章将为 Hom–张量及标量变换引入伴随的定义、单位与余单位；一般极限、伴随的系统理论和完整的 Yoneda 引理\footnote{米田信夫（假名：ヨネダ・ノブオ，罗马音：Nobuo Yoneda，1930-1996），日本数学家，提出 Yoneda 引理，并研究 Ext 乘积。}留到第四卷讨论。

% !TeX root = ../../Book2.tex
% 纲要标题：### 5.1 范畴与函子
% 纲要位置：计划纲要.md，第 952 行。

\section{范畴与函子}
\label{b2:sec:0501}

模的直和、商模和自由模虽然由不同公式构造，它们都通过同态说明自己的用途。范畴语言把“有哪些对象”与“对象之间允许哪些映射”一并记录下来。先要明确允许的映射：同一批集合配上不同的映射，可以得到不同的范畴，因而有不同的同构概念。章首问能否用从一个对象出发的映射来刻画它；比较这类映射性质时，还要知道改变来源或目标后，原有映射怎样随之变化。函子将记录这种变化，并保持映射之间的复合关系。

\ProseParagraphBreak
本书所说的范畴均为\term[jubuxiaofanchou]{局部小范畴}[locally small category]：任意两个对象之间的态射组成集合；对象全体则不一定组成集合。下面的定义采用这一约定。

% 纲要内容（仅作写作计划，尚非教材正文）：
%
% * 群、环、模和偏序的范畴
% * 协变、反变与对偶范畴
% * 忘却函子及自由函子
%

\subsection{群、环、模和偏序的范畴}
\label{b2:subsec:0501:01}

\begin{definition}\label{b2:def:category}
先给定一族对象。对每对对象 \(X,Y\)，给定一个集合，记为 \(\operatorname{Hom}_{\mathcal C}(X,Y)\)，也可以记为 \(\operatorname{Mor}_{\mathcal C}(X,Y)\)。其中的元素称为从 \(X\) 到 \(Y\) 的\term[taishe]{态射}[morphism]，写作 \(f:X\rightarrow Y\)，起点与终点也是态射数据的一部分。对态射 \(f:X\rightarrow Y\)、\(g:Y\rightarrow Z\) 给定复合 \(g\circ f:X\rightarrow Z\)，并要求：
\begin{enumerate}
\item 对可复合的三个态射，有 \(h\circ(g\circ f)=(h\circ g)\circ f\)。
\item 每个对象 \(X\) 有恒等态射 \(\operatorname{id}_X\)，满足 \(f\circ\operatorname{id}_X=f\) 及 \(\operatorname{id}_Y\circ f=f\)。
\end{enumerate}
满足这些条件的数据称为一个\term[fanchou]{范畴}[category] \(\mathcal C\)。
若态射 \(f:X\rightarrow Y\) 有态射 \(g:Y\rightarrow X\) 满足 \(gf=\operatorname{id}_X\)、\(fg=\operatorname{id}_Y\)，则称 \(f\) 为同构，写作 \(X\cong Y\)。
\end{definition}

\symidx{\(\operatorname{Hom}_{\mathcal C}(X,Y)\)：范畴中的态射集合}
\symidx{\(\operatorname{Mor}_{\mathcal C}(X,Y)\)：范畴中的态射集合的别称记号}
省略复合符号时，\(gf\) 总表示先做 \(f\)、后做 \(g\)。恒等态射唯一，因为两个候选 \(e,e'\) 满足 \(e=e e'=e'\)；逆态射也唯一，因为两个逆 \(g,h\) 满足 \(g=g(fh)=(gf)h=h\)。态射不必是集合之间的函数；定义只要求上述复合法则。

\ProseParagraphBreak
集合及全部函数构成 \(\mathbf{Set}\)，群及群同态构成 \(\mathbf{Grp}\)，交换群及群同态构成 \(\mathbf{Ab}\)。环及保单位元的环同态构成 \(\mathbf{Ring}\)；本书允许零环，因而从任意环到零环有唯一环同态。固定环 \(R\)，左 \(R\) 模及 \(R\) 线性映射构成 \(R\text{-}\mathbf{Mod}\)，右模范畴写作 \(\mathbf{Mod}\text{-}R\)。它们的复合来自函数复合，结合律和恒等律随之成立。在群、环和模的这些范畴中，同构正是已有的代数同构。

\symidx{\(R\text{-}\mathbf{Mod}\)：左模范畴}
\symidx{\(\mathbf{Mod}\text{-}R\)：右模范畴}
\symidx{\(\mathbf{Set}\)：集合范畴}
\symidx{\(\mathbf{Grp}\)：群范畴}
\symidx{\(\mathbf{Ab}\)：交换群范畴}
\symidx{\(\mathbf{Ring}\)：保单位元同态组成的环范畴}
对偏序集 \((P,\le)\)，把元素当作对象，规定 \(x\le y\) 时恰有一个态射 \(x\rightarrow y\)，否则没有态射。自反性提供恒等态射，传递性提供复合；所有同起点、同终点的态射至多一个，所以结合律自动成立。偏序中的两个对象同构当且仅当相等，因为双向态射意味着 \(x\le y\) 且 \(y\le x\)。与此不同，一个群 \(G\) 可以看成只有一个对象的范畴，该对象的全部态射为 \(G\)，复合为群乘法，每个态射都可逆。这里偏序的态射是由 \(x\le y\) 给出的形式箭头，单对象群的态射则是群元素；两者都直接规定箭头及其复合，不把态射定义为底集合之间的函数。

\ProseParagraphBreak
若范畴的对象也组成集合，就称它为\term[xiaofanchou]{小范畴}[small category]。本章涉及对象类时沿用固定的集合论宇宙约定，积与余积均由集合编号。

\subsection{协变、反变与对偶范畴}
\label{b2:subsec:0501:02}

\begin{definition}\label{b2:def:functor}
设 \(\mathcal C,\mathcal D\) 是范畴。给定赋值 \(F\)：将 \(\mathcal C\) 的每个对象 \(X\) 送到 \(\mathcal D\) 的对象 \(F(X)\)，将每个态射 \(f:X\rightarrow Y\) 送到 \(F(f):F(X)\rightarrow F(Y)\)。如果
\[
F(\operatorname{id}_X)=\operatorname{id}_{F(X)},\qquad F(gf)=F(g)F(f).
\]
就称 \(F\) 为从 \(\mathcal C\) 到 \(\mathcal D\) 的\term[hanzi]{函子}[functor]，也称\term[xiebian-hanzi]{协变函子}[covariant functor]。若赋值改为 \(F(f):F(Y)\rightarrow F(X)\)，并满足 \(F(\operatorname{id}_X)=\operatorname{id}_{F(X)}\)、\(F(gf)=F(f)F(g)\)，就称 \(F\) 为\term[fanbian-hanzi]{反变函子}[contravariant functor]。
\end{definition}

对可复合的 \(f:X\to Y\)、\(g:Y\to Z\)，协变和反变函子分别把复合关系写成左、右两个交换三角形。右图同时反转了箭头方向和复合次序。
\[
\begin{tikzcd}[column sep=2em]
F(X)\arrow[r,"F(f)"]\arrow[rr,bend right=25,"F(gf)"']&F(Y)\arrow[r,"F(g)"]&F(Z)
\end{tikzcd}
\qquad
\begin{tikzcd}[column sep=2em]
F(X)&F(Y)\arrow[l,"F(f)"']&F(Z)\arrow[l,"F(g)"']\arrow[ll,bend left=25,"F(gf)"]
\end{tikzcd}
\]

\begin{definition}\label{b2:def:opposite-category}
设 \(\mathcal C\) 是范畴。取与 \(\mathcal C\) 相同的对象，令
\[
\operatorname{Hom}_{\mathcal C^{\mathrm{op}}}(X,Y)=\operatorname{Hom}_{\mathcal C}(Y,X),
\]
复合由 \(\mathcal C\) 中反序的复合给出：把 \(f:X\rightarrow Y\) 记成反向态射 \(f^{\mathrm{op}}:Y\rightarrow X\)，规定 \(f^{\mathrm{op}}g^{\mathrm{op}}=(gf)^{\mathrm{op}}\)。所得范畴称为 \(\mathcal C\) 的\term[duiou-fanchou]{对偶范畴}[opposite category] \(\mathcal C^{\mathrm{op}}\)。
\end{definition}

\symidx{\(\mathcal C^{\mathrm{op}}\)：对偶范畴}
反转复合次序仍保留结合律与恒等律，所以这确实是范畴。从 \(\mathcal C\) 到 \(\mathcal D\) 的反变函子，就是从 \(\mathcal C^{\mathrm{op}}\) 到 \(\mathcal D\) 的普通函子：源范畴已经把箭头及复合次序反转，函子本身仍按协变公理保持复合。以后可以统一按这一方式理解反变函子。这里改变的是所有箭头方向；环的对合或某个向量空间的对偶是另外的具体构造，不能仅凭“对偶”二字混为一谈。

\ProseParagraphBreak
固定对象 \(A\)，赋值 \(X\mapsto\operatorname{Hom}_{\mathcal C}(A,X)\) 是到 \(\mathbf{Set}\) 的协变函子：\(f:X\rightarrow Y\) 将 \(u:A\rightarrow X\) 送到 \(fu\)，这是沿 \(f\) 改变目标时的后复合。相反，\(X\mapsto\operatorname{Hom}_{\mathcal C}(X,A)\) 为反变函子，\(f\) 将 \(v:Y\rightarrow A\) 送到 \(vf\)，这是沿 \(f\) 改变来源时的预复合。因此 Hom 对第一个变量反变，对第二个变量协变。复合次序可逐项核对：\(g(fu)=(gf)u\)，而 \((vg)f=v(gf)\)。这两类函子分别记录从固定对象出发和到固定对象为止的全部映射，以及这些映射随另一对象变化的方式。

\ProseParagraphBreak
在模范畴中，这些 Hom 集合还是交换群，相应映射保持加法，因此也可把目标取成 \(\mathbf{Ab}\)。一般非交换环上，不能把 \(\operatorname{Hom}_R(M,N)\) 无条件当作左 \(R\) 模；\cref{b2:prop:hom-central-action}只保证中心的标量作用。若 \(R=k\) 为域，\(V^*=\operatorname{Hom}_k(V,k)\) 与 \(f^*(\lambda)=\lambda f\) 给出向量空间的反变对偶函子，故 \((gf)^*=f^*g^*\)。

\subsection{忘却函子及自由函子}
\label{b2:subsec:0501:03}

把左模送到其底集合、把模同态送到同一个函数，得到\term[wangque-hanzi]{忘却函子}[forgetful functor]
\[
U:R\text{-}\mathbf{Mod}\longrightarrow\mathbf{Set}.
\]
忘却结构没有改变复合和恒等函数，故函子公理成立。当 \(R\ne0\) 时，\(U\) 的态射作用并不覆盖任意集合映射：正则模 \(R\) 的底集合上，平移 \(r\mapsto r+1_R\) 不保持零元，因而不是模同态。若 \(R=0\)，幺性模只有单元素的零模，上述反例便不适用。

\ProseParagraphBreak
对集合 \(X\)，取带指定基 \((e_x)_{x\in X}\) 的自由模 \(F(X)=R^{(X)}\)。对函数 \(a:X\rightarrow Y\)，\cref{b2:thm:free-module-universal}给出唯一模同态
\[
F(a):F(X)\longrightarrow F(Y),\qquad e_x\longmapsto e_{a(x)}.
\]
在每个基向量上，\(F(\operatorname{id}_X)\) 等于恒等映射，\(F(ba)\) 等于 \(F(b)F(a)\)；自由模泛性质的唯一性使这些等式在全模成立。因此 \(F:\mathbf{Set}\rightarrow R\text{-}\mathbf{Mod}\) 是\term[ziyou-hanzi]{自由函子}[free functor]。这里不为每个抽象模任意选基，而是从集合直接构造带有指定基的自由模。

\begin{proposition}\label{b2:prop:free-forgetful-bijection}
设 \(R\) 是环，\(X\) 是集合，\(M\) 是左 \(R\) 模。记 \(F(X)=R^{(X)}\)，其指定基为 \((e_x)_{x\in X}\)；记 \(U(M)\) 为 \(M\) 的底集合。限制到指定基给出双射
\begin{equation}\label{b2:eq:free-forgetful-bijection}
\begin{aligned}
\Phi_{X,M}:\operatorname{Hom}_R\bigl(F(X),M\bigr)
&\longrightarrow\operatorname{Hom}_{\mathbf{Set}}\bigl(X,U(M)\bigr),\\
h&\longmapsto\bigl(x\mapsto h(e_x)\bigr).
\end{aligned}
\end{equation}
若 \(X'\) 是集合、\(M'\) 是左 \(R\) 模、\(a:X'\rightarrow X\) 为函数、\(b:M\rightarrow M'\) 为模同态，取任意 \(h\in\operatorname{Hom}_R(F(X),M)\)。沿 \(a\) 改变生成元集合时，模同态变为 \(hF(a)\)，指定基像的函数则预复合 \(a\)；沿 \(b\) 改变目标模时，模同态变为 \(bh\)，指定基像的函数则后复合底集合映射 \(U(b)\)。同时作这两种改变，便有
\[
\Phi_{X',M'}\bigl(bhF(a)\bigr)=U(b)\,\Phi_{X,M}(h)\,a.
\]
\end{proposition}
\begin{proof}
给定函数 \(v:X\rightarrow U(M)\)，\cref{b2:thm:free-module-universal}说明存在唯一同态将每个 \(e_x\) 送到 \(v(x)\)，这证明双射性。对 \(x'\in X'\)，上式左端的值为 \(b\Bigl(h\bigl(e_{a(x')}\bigr)\Bigr)\)，右端也是同一个元素，故相等。
\end{proof}

该相容式说明，这一双射在改变生成元集合和目标模时都保持一致。下一节把这种“一族映射同时与所有态射相容”的条件称为自然性；张量积与 Hom 的对应也将满足同类条件。

% !TeX root = ../../Book2.tex
% 纲要标题：### 5.2 自然性
% 纲要位置：计划纲要.md，第 958 行。

\section{自然性}
\label{b2:sec:0502}

章首用自由模和商模的映射性质刻画对象；当输入对象改变时，这些映射对应也要与复合相容。上一节的自由--忘却双射已经显示了这种要求。“不依赖基”说明一个公式不需要先选坐标，却还没有完整说明不同对象之间的关系。自然性要求：先把对象上的向量或态射送到另一个对象，再使用构造，与先使用构造、再施加诱导映射，结果相同。第一章的双对偶求值映射也满足这种条件；现在把它作为一族映射来讨论。

% 纲要内容（仅作写作计划，尚非教材正文）：
%
% * 自然变换与自然同构
% * 恒等函子的自然自变换与环中心
% * 与选基依赖的同构比较
% * 模同构和函子同构的不同层次
% * 挠商的逐对象分裂与自然分裂
%

\subsection{自然变换与自然同构}
\label{b2:subsec:0502:01}

\begin{definition}\label{b2:def:natural-transformation}
设 \(\mathcal C,\mathcal D\) 是范畴，\(F,G:\mathcal C\rightarrow\mathcal D\) 是两个函子。给 \(\mathcal C\) 的每个对象 \(X\) 指定一个 \(\mathcal D\) 中的态射 \(\eta_X:F(X)\rightarrow G(X)\)。若对每个 \(f:X\rightarrow Y\) 都有
\begin{equation}\label{b2:eq:naturality-square}
G(f)\eta_X=\eta_YF(f).
\end{equation}
就称这族态射为\term[ziran-bianhuan]{自然变换}[natural transformation] \(\eta:F\Rightarrow G\)，并称 \(\eta_X\) 为其在 \(X\) 处的分量。若有自然变换 \(\theta:G\Rightarrow F\) 使每个 \(\theta_X\eta_X\) 和 \(\eta_X\theta_X\) 都是恒等态射，则称 \(\eta\) 为\term[ziran-tonggou]{自然同构}[natural isomorphism]。
\end{definition}

式\eqref{b2:eq:naturality-square}也可画成交换图：
\[
\begin{tikzcd}
F(X)\arrow[r,"\eta_X"]\arrow[d,"F(f)"']&G(X)\arrow[d,"G(f)"]\\
F(Y)\arrow[r,"\eta_Y"']&G(Y).
\end{tikzcd}
\]
从 \(F(X)\) 出发，先用 \(\eta_X\) 再用 \(G(f)\)，与先用 \(F(f)\) 再用 \(\eta_Y\)，必须得到相同的态射。因而不同对象处的分量不能任意选择：每个 \(f\) 都把 \(\eta_X\) 与 \(\eta_Y\) 联系起来。定义中必须给出两个函子以及它们对态射的作用；同一个对象公式若配上不同的态射作用，自然性条件就可能不同。

\ProseParagraphBreak
两个反变函子之间的自然变换，按它们从 \(\mathcal C^{\mathrm{op}}\) 出发来理解。原态射 \(f:X\rightarrow Y\) 在反范畴中指向相反方向，因此对应的方块为
\[
\begin{tikzcd}
F(Y)\arrow[r,"\eta_Y"]\arrow[d,"F(f)"']&G(Y)\arrow[d,"G(f)"]\\
F(X)\arrow[r,"\eta_X"']&G(X).
\end{tikzcd}
\]
所以相容式是 \(G(f)\eta_Y=\eta_XF(f)\)，两个水平分量的下标也随箭头方向一起交换。

\begin{proposition}\label{b2:prop:natural-isomorphism-components}
设 \(\mathcal C,\mathcal D\) 是范畴，\(F,G,H:\mathcal C\to\mathcal D\) 是函子。自然变换 \(\eta:F\Rightarrow G\) 是自然同构，当且仅当每个 \(\eta_X\) 都是 \(\mathcal D\) 中的同构。自然变换可以按分量复合；若 \(\theta:G\Rightarrow H\)，其与 \(\eta\) 的复合分量为 \((\theta\eta)_X=\theta_X\eta_X\)。
\end{proposition}
\begin{proof}
有自然逆时，每个分量当然有逆。反过来，若分量都可逆，由 \(G(f)\eta_X=\eta_YF(f)\) 两边分别复合逆映射，得到
\[
F(f)\eta_X^{-1}=\eta_Y^{-1}G(f).
\]
所以分量逆组成自然变换 \(G\Rightarrow F\)，并在每个对象上与 \(\eta\) 互逆。

对复合，由两个自然性等式有
\[
H(f)\theta_X\eta_X=\theta_YG(f)\eta_X=\theta_Y\eta_YF(f),
\]
故它仍自然。复合的结合律与恒等变换都逐分量继承自 \(\mathcal D\)。
\end{proof}

例如在交换群范畴上，固定整数 \(n\)，映射 \(\eta_A:A\rightarrow A\)、\(a\mapsto na\) 构成恒等函子到自身的自然变换，因为每个群同态 \(f:A\rightarrow B\) 都满足 \(f(na)=n\cdot f(a)\)。\(n=1,-1\) 时为自然同构，\(n=2\) 时在 \(\mathbb Z\) 上不满，在 \(\mathbb Z/2\mathbb Z\) 上也不单，因此不是自然同构。自然性要求分量与所有同态相容，分量本身可以不可逆。

\ProseParagraphBreak
任意环上的左模也可以问同一个问题：哪些操作能在所有模上同时定义，并与每个模同态相容？单个模有许多自同态，但这种同时相容的要求只留下同一个中心元素的作用。

\begin{proposition}\label{b2:prop:natural-center}
设 \(R\) 是含幺环，\(\mathcal C=R\text{-}\mathbf{Mod}\) 是幺性左 \(R\) 模及模同态的范畴，\(Z(R)=\{c\in R:cr=rc\text{ 对所有 }r\in R\}\)。恒等函子的自然自变换恰为
\[
\eta^c_M(m)=cm\qquad(M\in\mathcal C,\ m\in M,\ c\in Z(R)).
\]
逐分量相加和复合使这些自然自变换组成一个环，映射 \(c\mapsto\eta^c\) 是 \(Z(R)\) 到这个环的同构。自然自同构恰对应 \(Z(R)\) 的单位。
\end{proposition}
\begin{proof}
设 \(\eta\) 是自然自变换，令 \(c=\eta_R(1)\)。分量 \(\eta_R\) 左 \(R\) 线性，故 \(\eta_R(r)=rc\)。对任意左模 \(M\) 和 \(m\in M\)，同态 \(f_m:R\to M\)、\(r\mapsto rm\) 把 \(1\) 送到 \(m\)。将自然性等式 \(\eta_M f_m=f_m\eta_R\) 作用于 \(1\)，得到
\[
\eta_M(m)=f_m(c)=cm.
\]
尤其令 \(M=R,m=r\)，得到 \(\eta_R(r)=cr\)。与先前的 \(rc\) 比较，可知 \(c\in Z(R)\)。因此所有分量都由这一个元素确定。

反过来，若 \(c\in Z(R)\)，映射 \(m\mapsto cm\) 保持加法，且 \(c(rm)=r(cm)\)，故左 \(R\) 线性。每个模同态 \(f:M\to N\) 满足 \(f(cm)=cf(m)\)，所以这族映射自然。相加与复合分别对应 \(c+d\)、\(cd\)，恒等变换对应 \(1\)，而在正则模的 \(1\) 处取值恢复 \(c\)，从而得到环同构。若 \(\eta^c\) 可逆，其自然逆由某个 \(d\in Z(R)\) 给出，且 \(cd=dc=1\)；反向也成立。
\end{proof}

\subsection{自然同构与选基同构}
\label{b2:subsec:0502:02}

记 \(\mathbf{Vect}_k\) 为 \(k\) 向量空间及线性映射的范畴。两次应用对偶反转两次箭头，因而 \(D^2(V)=V^{**}\)、\(D^2(f)=f^{**}\) 是协变函子。求值映射
\[
j_V:V\longrightarrow V^{**},\qquad j_V(v)(\lambda)=\lambda(v)
\]
与恒等函子有相同方向。

\begin{proposition}\label{b2:prop:natural-double-dual}
设 \(k\) 是域，\(\mathbf{Vect}_k\) 是 \(k\) 向量空间及其线性映射的范畴，\(D^2(V)=V^{**}\)、\(D^2(f)=f^{**}\)。求值映射 \(j_V(v)(\lambda)=\lambda(v)\) 组成自然变换 \(j:\operatorname{Id}_{\mathbf{Vect}_k}\Rightarrow D^2\)。限制在有限维向量空间范畴上，它是自然同构。
\end{proposition}
\begin{proof}
设 \(f:V\rightarrow W\)，对 \(v\in V\)、\(\mu\in W^*\)，有
\[
\bigl(f^{**}j_V(v)\bigr)(\mu)
=j_V(v)(\mu f)=\mu\bigl(f(v)\bigr)
=j_W\bigl(f(v)\bigr)(\mu).
\]
因此 \(f^{**}j_V=j_Wf\)，这就是自然性。有限维时，\cref{b2:thm:double-dual}说明每个 \(j_V\) 为同构；由\cref{b2:prop:natural-isomorphism-components}，它们组成自然同构。
\end{proof}

这里对任意维数都验证了自然性，有限维只用于证明分量满射。无限维时则不保证满射：\cref{b2:ex:infinite-double-dual}对 \(V=k^{(\mathbb N)}\) 给出了一个不在求值映射像中的双对偶元素，因而不能把自然变换自动升级为自然同构。

\ProseParagraphBreak
前一节的 Hom 函子、对偶函子与这里的双对偶函子，对箭头的作用可以并列核对。表中前三行取 \(f:X\to Y\)，后两行取线性映射 \(f:V\to W\)。
\begin{table}[htbp]
\centering
\caption{常见函子的箭头方向}\label{b2:tab:functor-arrow-directions}
\begin{tabularx}{\linewidth}{@{}p{0.34\linewidth}X p{0.12\linewidth}@{}}
\toprule
构造 & 在输入箭头上的作用 & 类型 \\
\midrule
\(F:\mathcal C\to\mathcal D\) & \(F(f):F(X)\to F(Y)\) & 协变 \\
\(\operatorname{Hom}_{\mathcal C}(A,-)\) & \(u:A\to X\mapsto fu:A\to Y\) & 协变 \\
\(\operatorname{Hom}_{\mathcal C}(-,A)\) & \(v:Y\to A\mapsto vf:X\to A\) & 反变 \\
\(D(V)=V^*\) & \(f^*:W^*\to V^*\) & 反变 \\
\(D^2(V)=V^{**}\) & \(f^{**}:V^{**}\to W^{**}\) & 协变 \\
\bottomrule
\end{tabularx}
\end{table}
选基给出的 \(V\cong V^*\) 情况不同。在实一维空间 \(V=\mathbb R e\) 中，用基 \(e\) 规定 \(b_e(e)=e^*\)。改用 \(e'=2e\)，其对偶基为 \((e')^*=\dfrac12e^*\)，所以
\[
b_{e'}(e)=\dfrac12b_{e'}(e')=\dfrac14e^*,
\]
与 \(b_e(e)=e^*\) 不同。两者都是同构，但并非同一个映射。这只说明选基构造给出的同构不典范，还没有排除另一种构造可能给出自然同构；要证明不存在自然同构，须检验所有态射对分量施加的限制。还须留意，普通对偶函子是反变的，不能直接把 \(\operatorname{Id}\) 与它当作同一对范畴之间的两个协变函子来写自然同构。

\ProseParagraphBreak
现在只取有限维实向量空间为对象、线性同构为态射，得到一个群胚，即所有态射都可逆的范畴。在这个群胚上，\(V\mapsto V^*\)、\(f\mapsto(f^{-1})^*\) 构成协变函子；每个 \(V\) 都与 \(V^*\) 同构，但无法选出与所有线性同构相容的一族同构 \(b_V:V\rightarrow V^*\)。若有这样一族映射，则对任意线性同构 \(f:V\to W\)，自然性要求
\[
 (f^{-1})^*\circ b_V=b_W\circ f.
\]
同构 \(b_V:V\to V^*\) 等价于非退化双线性型 \(B_V(v,w)=b_V(v)(w)\)。上面的自然性等价于 \(B_W(fv,fw)=B_V(v,w)\)，即要求每个线性同构都保持这些型。对一个固定的非零空间，伸缩 \(2I\) 就应保持这个型；但双线性给出 \(B_V(2v,2w)=4B_V(v,w)\)，已经与非退化性冲突。将这一伸缩代回自然性公式，便得到
\[
\dfrac12 b_V=2b_V,
\]
于是 \(b_V=0\)，不可能是同构。若另给内积并只允许保持内积的同构，情形会改变；“是否自然”总要相对于明确的对象和态射来判断。

\subsection{模同构与函子同构}
\label{b2:subsec:0502:03}

设 \(F,G:\mathcal C\rightarrow R\text{-}\mathbf{Mod}\)。逐个证明 \(F(X)\cong G(X)\)，只说明每个对象处存在某个同构；为每个 \(X\) 选定同构 \(\eta_X\)，才得到一族具体分量；这族分量还须对每个态射满足式\eqref{b2:eq:naturality-square}，才能组成自然同构。任意选定的分量未必相容，甚至可能根本不存在相容的选择。

\begin{example}[逐对象同构而非自然同构]\label{b2:ex:conjugation-not-natural}
对复向量空间 \(V\)，令 \(\overline V\) 与 \(V\) 有相同的加法群，元素记为 \(\bar v\)，标量作用改为
\[
\lambda\bar v=\overline{\bar\lambda v}.
\]
对复线性映射 \(f:V\rightarrow W\)，令 \(C(f)(\bar v)=\overline{f(v)}\)，并令 \(C(V)=\overline V\)。这给出有限维复向量空间范畴上的函子 \(C\)。每个 \(V\) 都与 \(C(V)\) 同构，恒等函子却不与 \(C\) 自然同构。
\end{example}
\begin{proof}
\(C(f)\) 为复线性，因为
\[
C(f)(\lambda\bar v)=\overline{f(\bar\lambda v)}
=\overline{\bar\lambda\cdot f(v)}=\lambda\cdot C(f)(\bar v).
\]
它保持复合与恒等映射，故是函子。

先在 \(V=\mathbb Ce\) 中区分两种映射。取同名副本的集合映射 \(J:V\to\overline V\)、\(J(v)=\bar v\)，按共轭空间的标量作用满足
\[
J(ie)=\overline{ie}=-i\cdot\bar e=-iJ(e).
\]
另一映射 \(b:V\to\overline V\) 则指定 \(b(e)=\bar e\)，再作复线性延拓，因而
\[
b(ie)=i\cdot b(e)=i\cdot\bar e=\overline{-ie}.
\]
两映射在 \(e\) 上相同，却在 \(ie\) 上不同：\(J\) 是共轭线性的，而 \(b\) 是复线性的。一般地，若 \(e_1,\ldots,e_n\) 是 \(V\) 的基，指定 \(e_j\mapsto\bar e_j\) 并作复线性延拓就给出 \(V\cong\overline V\)。这个延拓将 \(\sum_j a_je_j\) 送到 \(\overline{\sum_j\bar a_je_j}\)，并非直接取同名元素。

要检测两函子对态射作用的差异，可取 \(f=i\operatorname{id}_V\)：恒等函子保留乘 \(i\) 的作用，共轭函子却将它变为乘 \(-i\)，因为 \(\overline{iv}=(-i)\bar v\)。若有自然变换 \(\eta:\operatorname{Id}\Rightarrow C\)，自然性就要求
\[
-i\eta_V=C(f)\eta_V=\eta_Vf=i\eta_V.
\]
复数域中特征为零，故 \(\eta_V=0\)。非零空间上的分量不可能可逆，因此不存在自然同构。
\end{proof}

本例中，每个对象处都有同构，选基也能给出一族具体同构，但乘 \(i\) 的态射排除了任何自然的选择。另一方面，选择本身并不必然破坏自然性：若选择后同时恰当地定义函子对态射的作用，也可能得到自然同构，\secref{b2:sec:0504}的矩阵范畴将给出这种例子。

\ProseParagraphBreak
前三章讨论过短正合列的分裂，第四章也证明有限生成交换群的挠商是自由群，因此每个挠商同态都有截面。但逐个选得截面，与让这些截面对所有群同态相容，仍是两个不同的问题。

\begin{proposition}\label{b2:prop:no-natural-torsion-splitting}
设 \(\mathbf{Ab}_{\mathrm{fg}}\) 是有限生成交换群及群同态的范畴，对对象 \(A\) 记挠子群为 \(T(A)\)，商群为 \(Q(A)=A/T(A)\)。每个同态 \(f:A\to B\) 诱导
\[
Q(f):Q(A)\to Q(B),\qquad a+T(A)\longmapsto f(a)+T(B),
\]
这些映射使 \(Q\) 成为函子，商同态 \(q_A:A\to Q(A)\) 组成自然变换 \(q:\operatorname{Id}\Rightarrow Q\)。每个 \(q_A\) 都有群同态截面，但不存在满足 \(q_As_A=\operatorname{id}_{Q(A)}\) 的自然变换 \(s:Q\Rightarrow\operatorname{Id}\)。
\end{proposition}
\begin{proof}
若 \(na=0\)，则 \(nf(a)=0\)，所以 \(f\) 将 \(T(A)\) 映入 \(T(B)\)。商映射的公式因而良定义，并逐陪集保持恒等和复合。等式 \(Q(f)q_A=q_Bf\) 给出 \(q\) 的自然性。\cref{b2:thm:pid-module-structure}应用于有限生成 \(\mathbb Z\) 模 \(A\)，说明 \(Q(A)\) 是有限秩自由交换群；由\cref{b2:cor:free-quotient-splits}，每个 \(q_A\) 有截面。

假设存在这样的自然截面族，考虑非零商同态 \(f:\mathbb Z\to\mathbb Z/2\mathbb Z\)。在标准识别下，\(Q(\mathbb Z)=\mathbb Z\)、\(q_{\mathbb Z}=\operatorname{id}\)，故截面只能是 \(s_{\mathbb Z}=\operatorname{id}\)。另一方面，\(Q(\mathbb Z/2\mathbb Z)=0\)，所以 \(s_{\mathbb Z/2\mathbb Z}Q(f)=0\)。自然性却要求
\[
fs_{\mathbb Z}=s_{\mathbb Z/2\mathbb Z}Q(f),
\]
左边是 \(f\ne0\)，矛盾。因此逐对象的分裂不能提升为自然分裂。
\end{proof}

% !TeX root = ../../Book2.tex
% 纲要标题：### 5.3 泛性质与普遍构造
% 纲要位置：工作记录/计划纲要.md。

\section{泛性质与普遍构造}
\label{b2:sec:0503}

同一个商模可以用不同生成元和关系表示，同一个直和也可以有不同的集合实现。若要比较这些实现，最有用的问题是：哪些映射能唯一地经过它？本节把这类泛性质写成统一的语言，并亲自构造模的拉回与推出。所有环仍保单位元，所有模均为左模；这些构造不要求环交换。

% 纲要内容（仅作写作计划，尚非教材正文）：
%
% * 初始对象、终对象、积与余积
% * 无限直积与指定嵌入不满足余积唯一性
% * 模的拉回与推出
% * 由 Hom 映射性质辨认对象
% * 正则模作为同构的测试对象
%

\subsection{初始对象、终对象、积与余积}
\label{b2:subsec:0503:01}

\begin{definition}\label{b2:def:initial-terminal-object}
设 \(\mathcal C\) 是范畴。若对象 \(I\) 到 \(\mathcal C\) 的每个对象 \(X\) 恰有一个态射，就称 \(I\) 为\term[chushi-duixiang]{初始对象}[initial object]。若每个 \(X\) 到对象 \(T\) 恰有一个态射，就称 \(T\) 为\term[zhong-duixiang]{终对象}[terminal object]。若一个对象兼为初始对象和终对象，就称它为\term[ling-duixiang]{零对象}[zero object]。
\end{definition}

集合范畴的初始对象为空集，终对象为任一单点集，故没有零对象；模范畴中的零模则为零对象。保单位元同态的环范畴中，\(\mathbb Z\) 是初始对象：到 \(R\) 的映射必为 \(n\mapsto n1_R\)，而此公式确实保持环运算与单位元。零环是终对象，却不是初始对象，因为从零环到非零环不存在保单位元同态。

\ProseParagraphBreak
两个初始对象 \(I,I'\) 之间的唯一态射 \(u:I\rightarrow I'\)、\(v:I'\rightarrow I\) 满足 \(vu=\operatorname{id}_I\)、\(uv=\operatorname{id}_{I'}\)，因为相应自态射集合各只有一个元素。因此初始对象之间有唯一同构；反转箭头即得终对象的相同结论。零对象的存在还为每对对象给出零态射，即经过零对象的唯一复合；在模范畴中它就是通常的零映射。

\begin{definition}\label{b2:def:product-coproduct}
设 \(\mathcal C\) 是范畴，\((X_i)_{i\in I}\) 为由集合 \(I\) 编号的 \(\mathcal C\) 中对象族。若对象 \(P\) 及态射 \(p_i:P\rightarrow X_i\) 满足：对任意对象 \(T\) 及态射族 \(a_i:T\rightarrow X_i\)，存在唯一态射 \(a:T\rightarrow P\) 使全部 \(p_i a=a_i\)，就称它们为该对象族的\term[ji]{积}[product]。

若对象 \(Q\) 及态射 \(j_i:X_i\rightarrow Q\) 满足：对任意对象 \(T\) 及态射族 \(b_i:X_i\rightarrow T\)，存在唯一态射 \(b:Q\rightarrow T\) 使全部 \(bj_i=b_i\)，就称它们为该对象族的\term[yuji]{余积}[coproduct]。
\end{definition}

这里的积是 \((P,(p_i)_{i\in I})\)，余积是 \((Q,(j_i)_{i\in I})\)：泛性质刻画的是对象及其指定的结构映射。例如，把 \(k^2\) 作为两个 \(k\) 的积时，两个坐标投影也是数据的一部分；只给出 \(k^2\) 的同构类型，还没有指定这项积构造。

\ProseParagraphBreak
积与余积的分解方向相反。下面两个三角形分别对每个 \(i\in I\) 成立；虚线是同时满足全族相容条件的唯一态射。
\[
\begin{tikzcd}[column sep=3em,row sep=2.5em]
T\arrow[r,dashed,"\exists!\,a"]\arrow[dr,"a_i"']&P\arrow[d,"p_i"]\\
&X_i
\end{tikzcd}
\qquad
\begin{tikzcd}[column sep=3em,row sep=2.5em]
X_i\arrow[r,"j_i"]\arrow[dr,"b_i"']&Q\arrow[d,dashed,"\exists!\,b"]\\
&T
\end{tikzcd}
\]
在模范畴中，积由逐坐标运算的 \(\prod_i M_i\) 及坐标投影给出。给定 \(a_i:T\rightarrow M_i\)，唯一可能的映射是 \(a(t)=\bigl(a_i(t)\bigr)_i\)，它逐坐标保持加法和左标量乘法，故确为模同态。余积由有限支集直和 \(\bigoplus_iM_i\) 及坐标嵌入给出。给定 \(b_i:M_i\rightarrow T\)，定义
\[
b\bigl((m_i)_i\bigr)=\sum_i b_i(m_i).
\]
有限支集保证右边有意义；分配律保证线性；每个直和元素都是有限个 \(j_i(m_i)\) 之和，保证唯一性。这正是\cref{b2:thm:module-product-universal}与\cref{b2:thm:module-sum-universal}在范畴语言中的表达。

\ProseParagraphBreak
有限个模的积与余积可以由同一个模承载，但其指定映射和泛性质方向不同。无限时不能混用：例如对域 \(k\)，直和 \(k^{(\mathbb N)}\) 中每个元素只有有限个非零坐标，而积 \(k^{\mathbb N}\) 还包含 \((1,1,\ldots)\)。恒等映射族 \(k\rightarrow k\) 合成的映射 \(k\rightarrow k^{\mathbb N}\)、\(a\mapsto(a,a,\ldots)\) 就不能落入直和。空积为终对象，空余积为初始对象，均可直接从定义读出。

\begin{proposition}\label{b2:prop:infinite-product-not-coproduct}
设 \(k\) 是域，\(P=k^{\mathbb N}\)，并令 \(j_n:k\to P\) 将一个标量送到仅第 \(n\) 个坐标取该值的向量。则 \(\bigl(P,(j_n)_{n\in\mathbb N}\bigr)\) 不是 \(k\) 向量空间范畴中这族 \(k\) 的余积。
\end{proposition}
\begin{proof}
令 \(S=k^{(\mathbb N)}\subseteq P\)，取 \(w=(1,1,\ldots)\notin S\)。在子空间 \(S+kw\) 中，表达 \(s+aw\) 唯一：若 \(aw\in S\) 且 \(a\ne0\)，则 \(aw\) 有无限多个非零坐标，矛盾。因此
\[
\ell_0(s+aw)=a
\]
定义线性泛函。由\cref{b2:lem:functional-extension}，它可以延拓为 \(P\) 上的线性泛函 \(\ell\)，满足 \(\ell(w)=1\) 且 \(\ell|_S=0\)。每个 \(j_n(k)\) 都在 \(S\) 中，故 \(\ell j_n=0\) 对全部 \(n\) 成立。然而零映射与 \(\ell\) 是两个不同的映射 \(P\to k\)，却在所有 \(j_n\) 上给出同一族零映射，违反余积的唯一性。
\end{proof}

\begin{proposition}\label{b2:prop:universal-object-uniqueness}
设 \((X_i)_{i\in I}\) 是范畴 \(\mathcal C\) 中由集合 \(I\) 编号的对象族。同一对象族的任意两个积之间，存在唯一保持全部投影的同构；两个余积之间，存在唯一保持全部指定嵌入的同构。
\end{proposition}
\begin{proof}
设 \((P,p_i)\)、\((P',p'_i)\) 都为积。用 \(P\) 及其映射族 \((p_i)\) 代入 \(P'\) 的泛性质，得唯一 \(u:P\rightarrow P'\) 满足 \(p'_iu=p_i\)；交换角色得 \(v:P'\rightarrow P\)。于是 \(p_ivu=p_i\)，而 \(\operatorname{id}_P\) 也满足这些等式。\(P\) 的唯一性迫使 \(vu=\operatorname{id}_P\)，同理 \(uv=\operatorname{id}_{P'}\)。\(u\) 的唯一性已经包含在构造中。

对余积 \((Q,j_i)\)、\((Q',j'_i)\)，两次泛性质分别给出 \(u j_i=j'_i\)、\(v j'_i=j_i\)。其复合在每个 \(j_i\) 上等于恒等映射，由余积的唯一性得 \(vu=\operatorname{id}_Q\)、\(uv=\operatorname{id}_{Q'}\)。
\end{proof}

“唯一”修饰保持指定结构映射的同构；抽去这些映射后，对象本身可能有许多自同构。例如 \(k^2\) 有许多可逆线性变换，只有恒等映射同时保持两个坐标投影。

\subsection{模的拉回与推出}
\label{b2:subsec:0503:02}

若两个模 \(M,N\) 都映到同一个模 \(P\)，就可以寻找像相等的成对元素：哪些 \((m,n)\) 满足 \(f(m)=g(n)\)？同样，从任意测试对象出发的两条映射也应在 \(P\) 中具有相同复合。拉回把这些相容的映射对合成一条映射。反过来，若两条映射 \(u:P\to M\)、\(v:P\to N\) 有公共来源，推出要给出一个共同目标，使来自同一 \(p\in P\) 的 \(u(p)\) 与 \(v(p)\) 在那里相等。这两种要求在一般范畴中分别表述如下。

\begin{definition}\label{b2:def:pullback-pushout}
设 \(\mathcal C\) 是范畴。给定态射 \(f:M\rightarrow P\)、\(g:N\rightarrow P\)。若对象 \(L\) 及态射 \(p:L\rightarrow M\)、\(q:L\rightarrow N\) 满足 \(fp=gq\)，并且对每对满足 \(fa=gb\) 的态射 \(a:T\rightarrow M\)、\(b:T\rightarrow N\)，存在唯一 \(h:T\rightarrow L\) 满足 \(ph=a\)、\(qh=b\)，就称它们为 \(f,g\) 的\term[lahui]{拉回}[pullback]。

给定态射 \(u:P\rightarrow M\)、\(v:P\rightarrow N\)。若对象 \(Q\) 及态射 \(i:M\rightarrow Q\)、\(j:N\rightarrow Q\) 满足 \(iu=jv\)，并且对每对满足 \(au=bv\) 的态射 \(a:M\rightarrow T\)、\(b:N\rightarrow T\)，存在唯一 \(h:Q\rightarrow T\) 满足 \(hi=a\)、\(hj=b\)，就称它们为 \(u,v\) 的\term[tuichu]{推出}[pushout]。
\end{definition}

四种构造都由测试对象的映射决定，区别在于箭头的方向以及拉回、推出所需的相容等式。下表中的 \(T\) 可为范畴中的任意对象。
\begin{table}[htbp]
\centering
\caption[积、余积、拉回与推出]{积、余积、拉回与推出的测试箭头}\label{b2:tab:universal-constructions}
\begin{tabularx}{\linewidth}{@{}p{0.13\linewidth}p{0.27\linewidth}X@{}}
\toprule
构造 & 给定图形 & 测试箭头与唯一分解 \\
\midrule
积 & \((X_i)_{i\in I}\) & 各 \(a_i:T\to X_i\) 唯一汇成 \(T\to\prod_iX_i\)。 \\
余积 & \((X_i)_{i\in I}\) & 各 \(b_i:X_i\to T\) 唯一汇成 \(\coprod_iX_i\to T\)。 \\
拉回 & \(M\xrightarrow{f}P\xleftarrow{g}N\) & \(a:T\to M\)、\(b:T\to N\) 满足 \(fa=gb\) 时，唯一汇成 \(T\to M\times_PN\)。 \\
推出 & \(M\xleftarrow{u}P\xrightarrow{v}N\) & \(a:M\to T\)、\(b:N\to T\) 满足 \(au=bv\) 时，唯一汇成 \(M\amalg_PN\to T\)。 \\
\bottomrule
\end{tabularx}
\end{table}
积与拉回接收从测试对象出发的箭头；余积与推出则把指向测试对象的箭头合成一条。拉回和推出还要求这两条箭头在公共对象处相容。

\ProseParagraphBreak
在模范畴中，拉回可以在 \(M\oplus N\) 内选择满足 \(f(m)=g(n)\) 的元素。推出则从同一份直和开始，把公共来源的两个像识别起来：直和内的 \((u(p),0)\) 与 \((0,v(p))\) 之差为 \((u(p),-v(p))\)，所以应将这些差值视为零。映射
\[
d:P\longrightarrow M\oplus N,\qquad
p\longmapsto\bigl(u(p),-v(p)\bigr)
\]
是模同态，它的像已是子模。因此拉回由满足约束的子模给出，推出由加入这些关系的商模给出。

\begin{proposition}\label{b2:prop:module-pullback-pushout}
设 \(R\) 是环，\(f:M\to P\)、\(g:N\to P\) 及 \(u:P\to M\)、\(v:P\to N\) 是左 \(R\) 模同态。在左 \(R\) 模范畴中，这两对同态的拉回和推出分别可取为
\begin{align}
M\times_PN&=\bigl\{\,(m,n)\in M\oplus N\mid f(m)=g(n)\,\bigr\},\label{b2:eq:module-pullback}\\
M\amalg_PN&=(M\oplus N)/\Bigl\{\,\bigl(u(p),-v(p)\bigr)\mid p\in P\,\Bigr\}.\label{b2:eq:module-pushout}
\end{align}
拉回的映射为坐标投影，推出的映射为 \(i(m)=\bigl[(m,0)\bigr]\)、\(j(n)=\bigl[(0,n)\bigr]\)。
\end{proposition}
\begin{proof}
若 \(f(m)=g(n)\)、\(f(m')=g(n')\)，则
\[
f(m+m')=g(n+n'),\qquad f(rm)=g(rn),
\]
所以式\eqref{b2:eq:module-pullback}定义子模，两个投影为模同态并满足相容条件。给定 \(a,b\) 且 \(fa=gb\)，令 \(h(t)=\bigl(a(t),b(t)\bigr)\)。相容条件保证其像在拉回中，坐标公式保证线性；两个投影的值又决定元素的两个坐标，所以 \(h\) 唯一。

式\eqref{b2:eq:module-pushout}的分母是同态 \(d:P\rightarrow M\oplus N\)、\(p\mapsto\bigl(u(p),-v(p)\bigr)\) 的像，因而是子模。商模存在，且
\[
iu(p)-jv(p)=\Bigl[\bigl(u(p),-v(p)\bigr)\Bigr]=0.
\]
若 \(au=bv\)，定义 \(s:M\oplus N\rightarrow T\)、\(s(m,n)=a(m)+b(n)\)。它线性，并在 \(d(P)\) 上为零，所以
\[
h\Bigl(\bigl[(m,n)\bigr]\Bigr)=a(m)+b(n)
\]
不依赖代表元，给出所需模同态。每个商元素都有形式 \(i(m)+j(n)\)，任何满足 \(hi=a\)、\(hj=b\) 的同态都必须有此公式，因而唯一。
\end{proof}

这些定义可用一对相反方向的交换图比较。外侧实线为任意满足定义中相容条件的 \(a,b\)，虚线为唯一的分解：
\[
\begin{tikzcd}[column sep=2.1em,row sep=2em]
T\arrow[rd,dashed,"\exists!\,h"]\arrow[rdd,bend right=14,"a"']\arrow[rrd,bend left=14,"b"]&&\\
&M\times_PN\arrow[r,"q"]\arrow[d,"p"']&N\arrow[d,"g"]\\
&M\arrow[r,"f"']&P
\end{tikzcd}
\qquad
\begin{tikzcd}[column sep=2.1em,row sep=2em]
P\arrow[r,"v"]\arrow[d,"u"']&N\arrow[d,"j"]\arrow[ddr,bend left=14,"b"]&\\
M\arrow[r,"i"']\arrow[drr,bend right=14,"a"']&M\amalg_PN\arrow[dr,dashed,"\exists!\,h"]&\\
&&T.
\end{tikzcd}
\]
交换性只是条件的一半，还必须验证任意相容映射的唯一分解。用两次泛性质构造比较态射，并用唯一性检验两个复合为恒等，就能像\cref{b2:prop:universal-object-uniqueness}一样证明拉回、推出在保持图中映射的同构意义下唯一。

\ProseParagraphBreak
例如在整数模中，\(f:\mathbb Z\rightarrow\mathbb Z\) 为乘二，\(g\) 为乘三。拉回的条件是 \(2m=3n\)，故 \(m=3t,n=2t\)，它与 \(\mathbb Z\) 同构，但两个投影分别为乘三、乘二。把两映射改读为从公共定义域出发的 \(u=2\)、\(v=3\)，推出为
\[
\mathbb Z^2/\mathbb Z(2,-3)\cong\mathbb Z,
\qquad \bigl[(m,n)\bigr]\longmapsto3m+2n.
\]
该映射满射，因为 \(3-2=1\)；其核由 \(3m=-2n\) 确定为 \(\mathbb Z(2,-3)\)，故\cref{b2:thm:module-first-isomorphism}给出同构。此时两个结构映射仍分别是乘三、乘二。同构类型相同并不消除它们所记录的原映射。

\ProseParagraphBreak
取零模到 \(P\) 的零映射，拉回可由 \(\ker f\) 及核嵌入 \(\iota\) 给出；取 \(P\) 到零模的零映射，推出可由 \(\operatorname{coker}u=M/u(P)\) 及商映射 \(q\) 给出。相应的交换图为
\[
\begin{tikzcd}[column sep=3em,row sep=2.2em]
\ker f\arrow[r,"\iota"]\arrow[d]&M\arrow[d,"f"]\\
0\arrow[r]&P
\end{tikzcd}
\qquad
\begin{tikzcd}[column sep=3em,row sep=2.2em]
P\arrow[r,"u"]\arrow[d]&M\arrow[d,"q"]\\
0\arrow[r]&\operatorname{coker}u.
\end{tikzcd}
\]
左图把 \(f(m)=0\) 作为子模条件，右图把 \(u(P)\) 中的元素作为要模掉的关系。核与余核因此是拉回、推出的特殊情形，也与\cref{b2:def:module-kernel-image-cokernel}中的具体定义一致。

\subsection{由 Hom 映射性质辨认对象}
\label{b2:subsec:0503:03}

前面通过泛性质证明了积、余积等构造在保持指定映射的同构意义下唯一。把它们的映射条件写成 Hom 集合之间的自然对应，还可以从映射性质辨认对象。以积 \((P,p_i)\) 为例，其泛性质给出双射
\[
\Phi_T:\operatorname{Hom}_{\mathcal C}(T,P)\longrightarrow
\prod_i\operatorname{Hom}_{\mathcal C}(T,X_i),\qquad
h\longmapsto(p_ih)_i.
\]
对 \(a:T'\rightarrow T\)，两侧都通过预复合 \(a\) 改变输入，得到交换方块
\[
\begin{tikzcd}[column sep=4em,row sep=2.5em]
\operatorname{Hom}_{\mathcal C}(T,P)
\arrow[r,"\Phi_T"]\arrow[d,"-\circ a"']&
\prod_i\operatorname{Hom}_{\mathcal C}(T,X_i)\arrow[d,"-\circ a"]\\
\operatorname{Hom}_{\mathcal C}(T',P)\arrow[r,"\Phi_{T'}"']&
\prod_i\operatorname{Hom}_{\mathcal C}(T',X_i).
\end{tikzcd}
\]
右侧的预复合逐坐标进行，两条路径都把 \(h\) 送到 \((p_iha)_i\)，所以这些双射关于 \(T\) 自然。拉回则把右端换成满足 \(fa=gb\) 的映射对；推出的对应从 \(\operatorname{Hom}_{\mathcal C}(Q,T)\) 出发，关于 \(T\) 协变。它们不只数出了映射有多少，而是给出了与复合相容的具体对应。

\begin{proposition}[由映射检测同构]\label{b2:prop:hom-detects-isomorphism}
设 \(\mathcal C\) 是范畴，\(u:P\rightarrow Q\) 为 \(\mathcal C\) 中的态射。若对每个对象 \(T\)，后复合映射
\[
u_*:\operatorname{Hom}_{\mathcal C}(T,P)\longrightarrow
\operatorname{Hom}_{\mathcal C}(T,Q),\qquad h\longmapsto uh
\]
都是双射，则 \(u\) 为同构。更一般地，任何自然变换
\(\alpha:\operatorname{Hom}_{\mathcal C}(-,P)\Rightarrow
\operatorname{Hom}_{\mathcal C}(-,Q)\)
都唯一地由某个态射 \(u:P\rightarrow Q\) 的后复合给出；这个自然变换为自然同构，当且仅当 \(u\) 为同构。
\end{proposition}
\begin{proof}
先取 \(T=Q\)。满射性给出 \(v:Q\rightarrow P\)，使 \(uv=\operatorname{id}_Q\)。再取 \(T=P\)，有 \(u(vu)=u=u\operatorname{id}_P\)；这次的单射性给出 \(vu=\operatorname{id}_P\)，所以 \(u\) 可逆。

对自然变换 \(\alpha\)，若它由某个 \(u:P\to Q\) 后复合给出，在 \(T=P\)、\(h=\operatorname{id}_P\) 处就必有 \(\alpha_P(\operatorname{id}_P)=u\operatorname{id}_P=u\)。因此应当用恒等态射的像恢复候选态射，令 \(u=\alpha_P(\operatorname{id}_P)\)。给定 \(h:T\rightarrow P\)，用 \(h\) 的预复合及 \(\alpha\) 的自然性，得到
\[
\alpha_T(h)=\alpha_T(\operatorname{id}_P h)
=\alpha_P(\operatorname{id}_P)h=uh.
\]
所以 \(\alpha\) 的全部分量都由 \(u\) 后复合给出，在 \(T=P\) 的恒等态射处取值又说明 \(u\) 唯一。若 \(\alpha\) 是自然同构，其分量都为双射，第一部分说明 \(u\) 是同构。反过来，若 \(u\) 是同构，后复合 \(u^{-1}\) 就给出每个分量的逆，所以 \(\alpha\) 是自然同构。
\end{proof}

\begin{corollary}[正则模检测同构]\label{b2:cor:regular-module-detects}
设 \(R\) 是环，\(f:M\to N\) 是左 \(R\) 模同态。后复合映射
\[
\operatorname{Hom}_R(R,f):\operatorname{Hom}_R(R,M)
\longrightarrow\operatorname{Hom}_R(R,N)
\]
为双射，当且仅当 \(f\) 是模同构。
\end{corollary}
\begin{proof}
由\cref{b2:prop:regular-hom-endomorphism}，求值 \(\varphi\mapsto\varphi(1_R)\) 分别将两个 Hom 群识别为 \(M,N\)。而 \((f\varphi)(1_R)=f(\varphi(1_R))\)，所以后复合映射在这些求值对应下就是 \(f\) 本身。若它双射，\(f\) 是双射模同态，因而是\cref{b2:def:module-homomorphism}所定义的模同构；逆映射保持模运算。反过来，后复合 \(f^{-1}\) 给出 \(\operatorname{Hom}_R(R,f)\) 的逆。
\end{proof}

一般固定测试对象并没有正则模的这种能力。例如，包含同态 \(\mathbb Z\to\mathbb Q\) 不是同构，但它在 \(\operatorname{Hom}_{\mathbb Z}(\mathbb Z/2\mathbb Z,-)\) 下成为两个零群之间的双射：从 \(\mathbb Z/2\mathbb Z\) 出发的同态，其生成元像必须被 \(2\) 零化，而 \(\mathbb Z,\mathbb Q\) 中满足这一条件的都只有零。因此“对每个 \(T\)”不能无条件换成任意一个固定 \(T\)。

\ProseParagraphBreak
第四卷的 Yoneda 引理将把上述第二个 Hom 函子换成一般集合值反变函子，以同样的恒等态射求值思想描述自然变换。这里要求的是与所有复合相容的具体对应；即使各个 Hom 集合大小恰好相同，也没有自动得到双射的自然相容性。

% !TeX root = ../../Book2.tex
% 纲要标题：### 5.4 范畴等价
% 纲要位置：计划纲要.md，第 970 行。

\section{范畴等价}
\label{b2:sec:0504}

上一节用积等构造的映射性质来辨认对象；若要比较两个范畴，则还须比较每一对对象之间的态射及其复合。选择坐标把线性映射写成矩阵，往往没有改变我们要研究的结构，却改变了对象的具体外形。范畴等价使这种关系成为可证明的陈述：它要求对象可以在同构意义下互相对应，也要求全部态射及其复合得到保留。

% 纲要内容（仅作写作计划，尚非教材正文）：
%
% * 全、忠实与本质满
% * 范畴等价、同构与泛性质
% * 坐标选择所得拟逆之间的自然同构
% * 为第16章 Morita 理论准备最低限度语言
%
% ---
%

\subsection{全、忠实与本质满}
\label{b2:subsec:0504:01}

\begin{definition}\label{b2:def:fully-faithful-essentially-surjective}
设 \(\mathcal C,\mathcal D\) 是范畴，\(F:\mathcal C\rightarrow\mathcal D\) 是函子。对 \(\mathcal C\) 中每对对象 \(X,Y\)，\(F\) 给出函数
\[
F_{X,Y}:\operatorname{Hom}_{\mathcal C}(X,Y)\longrightarrow
\operatorname{Hom}_{\mathcal D}\bigl(F(X),F(Y)\bigr),\qquad f\longmapsto F(f).
\]
若这些函数全部满射，称 \(F\) 为\term[quan-hanzi]{全函子}[full functor]；若全部单射，称为\term[zhongshi-hanzi]{忠实函子}[faithful functor]；若全部双射，称为\term[quanzhongshi-hanzi]{全忠实函子}[fully faithful functor]。若 \(\mathcal D\) 的每个对象都同构于某个 \(F(X)\)，称 \(F\) \term[benzhi-man]{本质满}[essentially surjective]。
\end{definition}

这三个条件检查的位置不同：
\begin{table}[htbp]
\centering
\caption{全、忠实与本质满的核对位置}\label{b2:tab:functor-equivalence-conditions}
\begin{tabularx}{\linewidth}{@{}p{0.19\linewidth}X X@{}}
\toprule
条件 & 对每个 Hom 映射的要求 & 对目标范畴对象的要求 \\
\midrule
忠实 & 单射 & 无 \\
全 & 满射 & 无 \\
本质满 & 无 & 每个对象同构于某个像对象 \\
\bottomrule
\end{tabularx}
\end{table}

“忠实”不合并固定起点、终点之间的不同态射，但不要求不同对象有不同的像；“全”要求目标中像对象之间的态射没有遗漏；“本质满”则要求目标对象在同构意义下没有遗漏。模到集合的忘却函子忠实：两个模同态若为同一个函数，就确实相等。当 \(R\ne0\) 时它不全，因为正则模底集合上的 \(r\mapsto r+1_R\) 不保持零元；零环的模均为单元素模，此时忘却函子反而全忠实。交换群到群的包含函子全忠实，因为交换群之间的群同态正是交换群同态；它不本质满，因为非交换群不能同构于交换群。

\ProseParagraphBreak
对偏序集之间的函子 \(F:P\rightarrow Q\)，函子性恰好是保序性。它总忠实，因为起点、终点固定后至多一个态射；它全当且仅当 \(F(x)\le F(y)\) 蕴含 \(x\le y\)。所以“保序”不等于“全”：把两元素反链送到两元素链的对应位置，是忠实而不全的函子。

\begin{proposition}\label{b2:prop:fully-faithful-reflects-isomorphism}
设 \(F:\mathcal C\to\mathcal D\) 是范畴间的函子。\(F\) 把 \(\mathcal C\) 中的同构送到 \(\mathcal D\) 中的同构。若 \(F\) 全忠实，则对 \(\mathcal C\) 中的态射 \(f:X\to Y\)，\(F(f)\) 可逆蕴含 \(f\) 可逆；对对象 \(X,Y\)，\(F(X)\cong F(Y)\) 蕴含 \(X\cong Y\)。
\end{proposition}
\begin{proof}
若 \(g=f^{-1}\)，函子公理给出 \(F(g)F(f)=F(gf)=\operatorname{id}\) 及反序的同样等式，故 \(F(g)\) 为逆。

现设 \(F\) 全忠实。若 \(F(f)\) 的逆为 \(v:F(Y)\rightarrow F(X)\)，要证明 \(f\) 可逆，须先把 \(v\) 提升为候选逆，再把两条逆关系从像中反映回来。全性给出 \(g:Y\rightarrow X\) 使 \(F(g)=v\)，于是
\[
\begin{aligned}
F(gf)&=vF(f)=\operatorname{id}_{F(X)}=F(\operatorname{id}_X),\\
F(fg)&=F(f)v=\operatorname{id}_{F(Y)}=F(\operatorname{id}_Y).
\end{aligned}
\]
忠实性分别给出 \(gf=\operatorname{id}_X\)、\(fg=\operatorname{id}_Y\)，故 \(g=f^{-1}\)。若只给定 \(a:F(X)\rightarrow F(Y)\) 为同构，先用全性把 \(a\) 提升为 \(f\)，再用刚证明的结论即可。
\end{proof}

\subsection{范畴等价、同构与泛性质}
\label{b2:subsec:0504:02}

\begin{definition}\label{b2:def:category-equivalence}
设 \(\mathcal C,\mathcal D\) 是范畴，\(F:\mathcal C\rightarrow\mathcal D\) 是函子。若存在函子 \(G:\mathcal D\rightarrow\mathcal C\) 及自然同构
\[
\eta:\operatorname{Id}_{\mathcal C}\Rightarrow GF,\qquad
\varepsilon:FG\Rightarrow\operatorname{Id}_{\mathcal D},
\]
则称 \(F\) 为\term[fanchou-dengjia]{范畴等价}[equivalence of categories]，\(G\) 为其\term[zhunni]{拟逆}[quasi-inverse]，写作 \(\mathcal C\simeq\mathcal D\)。若可要求 \(GF=\operatorname{Id}_{\mathcal C}\)、\(FG=\operatorname{Id}_{\mathcal D}\) 作为函子严格相等，则称 \(F\) 为\term[fanchou-tonggou]{范畴同构}[isomorphism of categories]。
\end{definition}
\symidx{\(\mathcal C\simeq\mathcal D\)：范畴等价}

定义中的 \(\eta,\varepsilon\) 都是自然同构，因而也可以取它们的逆，用相反方向表达同一个比较。重要的是 \(GF\) 与 \(\operatorname{Id}_{\mathcal C}\)、\(FG\) 与 \(\operatorname{Id}_{\mathcal D}\) 分别自然同构。范畴同构在对象上给出真正的双射；等价则允许多个互相同构的对象合并到一个代表。比如范畴 \(\mathcal E\) 有两个对象 \(a,b\)，每对对象之间恰有一个态射，复合由唯一性确定。把它们都送到仅有一个对象及恒等态射的范畴 \(\mathbf 1\)，并反向选取 \(a\)，便给出等价：每个对象到 \(a\) 的唯一态射组成所需自然同构。但这两个范畴对象数不同，故不可能同构。

\ProseParagraphBreak
构造拟逆时，本质满使目标对象能找到同构的像对象；选定这些代表后，就要把目标中的态射经同构搬到两个像对象之间。全性保证搬来的态射有原像，忠实性保证原像唯一，并将恒等和复合的等式反映回原范畴。下面把这些要求分别落实到对象和态射的构造中。

\begin{theorem}[范畴等价的判据]\label{b2:thm:equivalence-criterion}
设 \(\mathcal C,\mathcal D\) 是范畴，\(F:\mathcal C\rightarrow\mathcal D\) 是函子。
\begin{enumerate}
\item 若 \(F\) 是范畴等价，则它全忠实且本质满。
\item 若 \(F\) 全忠实，并已为每个 \(Y\in\mathcal D\) 同时指定对象 \(G(Y)\in\mathcal C\) 及同构 \(\varepsilon_Y:F(G(Y))\to Y\)，则 \(F\) 是范畴等价。
\end{enumerate}
\end{theorem}
\begin{proof}
先设 \(G,\eta,\varepsilon\) 给出等价。对每个 \(Y\in\mathcal D\)，同构 \(\varepsilon_Y:F\bigl(G(Y)\bigr)\rightarrow Y\) 证明本质满。若 \(f,g:X\rightarrow X'\) 满足 \(F(f)=F(g)\)，\(\eta\) 的自然性给出
\[
\eta_{X'}f=GF(f)\eta_X=GF(g)\eta_X=\eta_{X'}g.
\]
消去同构 \(\eta_{X'}\)，得 \(f=g\)，所以 \(F\) 忠实。为了随后验证提升来的态射确实具有所需的 \(F\) 像，还要用到 \(G\) 的忠实性。若 \(u,v:Y\rightarrow Y'\) 且 \(G(u)=G(v)\)，\(\varepsilon\) 的自然性给出
\[
u\varepsilon_Y=\varepsilon_{Y'}FG(u)
=\varepsilon_{Y'}FG(v)=v\varepsilon_Y.
\]
消去 \(\varepsilon_Y\) 得 \(u=v\)，所以 \(G\) 也忠实。

给定 \(a:F(X)\rightarrow F(X')\)，先用 \(G\) 得到 \(G(a):GF(X)\rightarrow GF(X')\)，再经 \(\eta_X\) 与 \(\eta_{X'}^{-1}\) 搬回 \(X,X'\) 之间，得到候选态射
\[
f=\eta_{X'}^{-1}G(a)\eta_X:X\longrightarrow X'.
\]
由 \(\eta\) 对 \(f\) 的自然性，\(GF(f)\eta_X=\eta_{X'}f=G(a)\eta_X\)。消去 \(\eta_X\)，得 \(G\bigl(F(f)\bigr)=G(a)\)；再由 \(G\) 忠实得 \(F(f)=a\)。故 \(F\) 全。

反过来，设 \(F\) 全忠实，并已同时指定所述 \(G(Y)\) 与 \(\varepsilon_Y\)。对 \(a:Y\rightarrow Y'\)，要定义 \(G(a)\)，应先要求它在 \(F\) 下的像与这些代表同构相容，即下图交换：
\[
\begin{tikzcd}[column sep=3em,row sep=2em]
F(G(Y))\arrow[r,dashed,"F(G(a))"]\arrow[d,"\varepsilon_Y"']&
F(G(Y'))\arrow[d,"\varepsilon_{Y'}"]\\
Y\arrow[r,"a"']&Y'.
\end{tikzcd}
\]
因而中间态射必须先经 \(\varepsilon_Y\) 到 \(Y\)，再施加 \(a\)，最后经 \(\varepsilon_{Y'}^{-1}\) 回到 \(F(G(Y'))\)。全性使这条搬来的态射有原像，忠实性使原像唯一，所以存在唯一 \(G(a):G(Y)\rightarrow G(Y')\) 满足
\begin{equation}\label{b2:eq:quasi-inverse-on-arrows}
F\bigl(G(a)\bigr)=\varepsilon_{Y'}^{-1}a\varepsilon_Y.
\end{equation}
当 \(a=\operatorname{id}_Y\) 时，右边为恒等态射，忠实性给出 \(G(\operatorname{id}_Y)=\operatorname{id}_{G(Y)}\)。对 \(a:Y\rightarrow Y'\)、\(b:Y'\rightarrow Y''\)，有
\[
\begin{aligned}
F\bigl(G(b)G(a)\bigr)
&=\varepsilon_{Y''}^{-1}b\varepsilon_{Y'}\varepsilon_{Y'}^{-1}a\varepsilon_Y\\
&=\varepsilon_{Y''}^{-1}ba\varepsilon_Y=F\bigl(G(ba)\bigr).
\end{aligned}
\]
忠实性给出 \(G(b)G(a)=G(ba)\)，故 \(G\) 是函子。式\eqref{b2:eq:quasi-inverse-on-arrows}恰好说 \(\varepsilon\) 是自然同构。

还需把 \(X\) 与 \(GF(X)\) 比较。将 \(Y=F(X)\) 代入已选的代表同构，得到 \(\varepsilon_{F(X)}:F(GF(X))\rightarrow F(X)\)。它的逆方向恰好是 \(F(X)\rightarrow F(GF(X))\)，因此用全忠实性唯一地提升这个逆，得到 \(\eta_X:X\rightarrow GF(X)\)，满足
\[
F(\eta_X)=\varepsilon_{F(X)}^{-1}.
\]
右端可逆，\cref{b2:prop:fully-faithful-reflects-isomorphism}说明 \(\eta_X\) 可逆。对 \(f:X\rightarrow X'\)，式\eqref{b2:eq:quasi-inverse-on-arrows}用于 \(F(f)\) 后给出
\[
\begin{aligned}
F\bigl(GF(f)\eta_X\bigr)
&=\varepsilon_{F(X')}^{-1}F(f)\varepsilon_{F(X)}\varepsilon_{F(X)}^{-1}\\
&=\varepsilon_{F(X')}^{-1}F(f)
=F\bigl(\eta_{X'}f\bigr).
\end{aligned}
\]
由忠实性，\(GF(f)\eta_X=\eta_{X'}f\)，所以 \(\eta\) 自然，所需等价得证。

\end{proof}

若 \(\mathcal C,\mathcal D\) 都是小范畴，本质满给出的候选对 \((X,a)\)、\(a:F(X)\xrightarrow{\sim}Y\) 组成非空集合，并由集合 \(\operatorname{Ob}(\mathcal D)\) 编号；选择公理因此能同时选出定理第二项所需的代表。故对小范畴，等价恰好等于全忠实且本质满。对任意大范畴，逐个对象的存在性还不自动给出这种统一选择，因此定理将已指定的同构代表作为数据。若按\secref{b2:sec:0501}的大小约定，两个范畴在某个更大宇宙中成为小范畴，则可在该宇宙中作出选择。这项判据及拟逆构造可对照 Riehl\cite[Theorem 1.5.9, pp. 33--34]{Riehl2016CategoryTheory}。

\ProseParagraphBreak
设 \(k\) 是域，定义矩阵范畴 \(\mathbf{Mat}_k\)：对象为非负整数，\(n\) 到 \(m\) 的态射是 \(m\times n\) 矩阵，复合为矩阵乘法，恒等态射为单位矩阵；零尺寸按唯一的空矩阵处理。函子
\[
J:\mathbf{Mat}_k\longrightarrow\mathbf{Vect}_k^{\mathrm{fd}},\qquad
n\longmapsto k^n,\quad A\longmapsto(v\mapsto Av)
\]
\symidx{\(\mathbf{Mat}_k,\mathbf{Vect}_k^{\mathrm{fd}}\)：矩阵范畴与有限维向量空间范畴}
全忠实，因为标准基下每个线性映射恰由一个矩阵表示；它本质满，因为有限维空间有基，因而同构于 \(k^n\)。这里固定一个包含 \(k\) 的宇宙，只考虑其中的有限维向量空间；在更大的宇宙中，两范畴都小，因而可用上述判据得到等价。

\ProseParagraphBreak
具体选定每个有限维 \(V\) 的坐标同构 \(c_V:k^{\dim_k V}\rightarrow V\)，定义拟逆
\[
H:\mathbf{Vect}_k^{\mathrm{fd}}\longrightarrow\mathbf{Mat}_k,
\qquad V\longmapsto\dim_kV,
\]
并把 \(f:V\rightarrow W\) 送到 \(c_W^{-1}fc_V\) 的矩阵。复合时中间的坐标同构相消，给出 \(H(gf)=H(g)H(f)\)，且 \(H(\operatorname{id}_V)=I_{\dim_kV}\)，所以 \(H\) 是函子。这里 \(c_V\) 的源为 \(JH(V)\)，目标为 \(V\)，所以它们组成的比较是 \(c:JH\Rightarrow\operatorname{Id}_{\mathbf{Vect}_k^{\mathrm{fd}}}\)。等式 \(fc_V=c_WJH(f)\) 给出自然性方块：
\[
\begin{tikzcd}[column sep=3em,row sep=2em]
JH(V)\arrow[r,"c_V"]\arrow[d,"JH(f)"']&V\arrow[d,"f"]\\
JH(W)\arrow[r,"c_W"']&W.
\end{tikzcd}
\]
每个 \(c_V\) 都可逆，故\cref{b2:prop:natural-isomorphism-components}使 \(c\) 成为自然同构。在矩阵一侧，\(HJ(n)=n\)，而标准基下 \(c_{k^n}^{-1}\) 的矩阵 \(\eta_n:n\to HJ(n)\) 满足
\[
\begin{aligned}
HJ(A)\eta_n&=c_{k^m}^{-1}Ac_{k^n}c_{k^n}^{-1}\\
&=c_{k^m}^{-1}A=\eta_mA
\qquad(A:n\to m).
\end{aligned}
\]
因此 \((\eta_n)\) 给出自然同构 \(\operatorname{Id}_{\mathbf{Mat}_k}\Rightarrow HJ\)。拟逆 \(H\) 及这些比较依赖于坐标选择，\(J\) 则只使用标准基。维数只描述对象；\(H\) 还必须记录同态的矩阵，才能保持态射与复合。

\begin{proposition}\label{b2:prop:coordinate-choice-natural}
设 \(k\) 是域，在有限维 \(k\) 向量空间范畴 \(\mathbf{Vect}_k^{\mathrm{fd}}\) 中，为每个对象 \(V\) 指定两族坐标同构 \(c_V,d_V:k^{\dim_kV}\to V\)。令 \(H_c,H_d:\mathbf{Vect}_k^{\mathrm{fd}}\to\mathbf{Mat}_k\) 在对象上均取 \(V\mapsto\dim_kV\)，在态射 \(f:V\to W\) 上分别取 \(c_W^{-1}fc_V\)、\(d_W^{-1}fd_V\) 的矩阵。则标准基下 \(d_V^{-1}c_V\) 的矩阵 \(P_V\) 组成自然同构 \(P:H_c\Rightarrow H_d\)。
\end{proposition}
\begin{proof}
对每个 \(f:V\to W\)，矩阵乘法给出
\[
H_d(f)P_V=d_W^{-1}fd_Vd_V^{-1}c_V
=d_W^{-1}fc_V=P_WH_c(f).
\]
这就是 \(P\) 的自然性。每个 \(P_V\) 都可逆，其逆为 \(c_V^{-1}d_V\) 的矩阵；由\cref{b2:prop:natural-isomorphism-components}，\(P\) 是自然同构。因而不同坐标选择虽然可能产生不同的拟逆，它们由这些换基矩阵自然地联系起来。
\end{proof}

泛性质只规定相容态射的存在与唯一性。等价的全忠实性保留每个 Hom 集合及复合，本质满使任意目标测试对象都能换成同构的像对象，因此这些存在与唯一性要求可以在两个范畴之间传递。

\begin{proposition}\label{b2:prop:equivalence-preserves-universal-objects}
设 \(F:\mathcal C\to\mathcal D\) 是范畴等价。\(F\) 把 \(\mathcal C\) 的初始对象、终对象分别送到 \(\mathcal D\) 的初始对象、终对象。对任意由集合 \(I\) 索引的对象族 \((X_i)_{i\in I}\)，若其积或余积存在，\(F\) 还把该积或余积连同指定映射送到 \((F(X_i))_{i\in I}\) 的积或余积。
\end{proposition}
\begin{proof}
由\cref{b2:thm:equivalence-criterion}的正向结论，\(F\) 全忠实且本质满。若 \(I_0\) 是初始对象，对任意 \(Y\in\mathcal D\)，选同构 \(a:F(T)\rightarrow Y\)。全忠实性与后复合 \(a\) 给出双射
\[
\operatorname{Hom}_{\mathcal C}(I_0,T)
\longrightarrow\operatorname{Hom}_{\mathcal D}\bigl(F(I_0),Y\bigr).
\]
左边恰一个元素，故 \(F(I_0)\) 初始。若 \(Z\) 是终对象，预复合 \(a^{-1}\) 把 \(\operatorname{Hom}_{\mathcal C}(T,Z)\) 双射到 \(\operatorname{Hom}_{\mathcal D}\bigl(Y,F(Z)\bigr)\)，故 \(F(Z)\) 是终对象。

设 \((P,p_i)\) 为 \((X_i)_{i\in I}\) 的积。给定 \(b_i:Y\rightarrow F(X_i)\)，仍选 \(a:F(T)\rightarrow Y\)。全忠实性给出唯一 \(c_i:T\rightarrow X_i\) 使 \(F(c_i)=b_i a\)。积的泛性质给出唯一 \(c:T\rightarrow P\)，使 \(p_i c=c_i\)。于是 \(b=F(c)a^{-1}:Y\rightarrow F(P)\) 满足 \(F(p_i)b=b_i\)。若另一态射 \(b'\) 也满足，先把 \(b'a:F(T)\rightarrow F(P)\) 唯一提升为 \(c':T\rightarrow P\)；忠实性给出 \(p_ic'=c_i\)，积的唯一性给出 \(c'=c\)，所以 \(b'=b\)。

对余积 \((Q,j_i)\)，给定 \(b_i:F(X_i)\rightarrow Y\)，把 \(a^{-1}b_i\) 唯一提升为 \(c_i:X_i\rightarrow T\)，由余积得到 \(c:Q\rightarrow T\)，再令 \(b=aF(c)\)。若 \(b'F(j_i)=b_i\)，把 \(a^{-1}b'\) 提升后，用余积的唯一性得到同一个 \(c\)。故余积也被保留。
\end{proof}

这里保留积的结论，是 \(F(P)\) 连同 \(F(p_i)\) 满足目标范畴的积泛性质，因而与任何另一个积存在唯一相容同构；不要求 \(F(P)\) 与某个底集合的坐标直积严格相等。余积也必须连同指定嵌入一起比较，只有对象的抽象同构类型不足以检验这种相容性。

\subsection{Morita 理论的范畴语言}
\label{b2:subsec:0504:03}

若两个环 \(R,S\) 同构，它们的模论只相差标量的名称。设 \(\varphi:R\rightarrow S\) 为环同构，对左 \(R\) 模 \(M\) 定义左 \(S\) 作用
\[
s\cdot m=\varphi^{-1}(s)m.
\]
在模同态上保持原映射，得到 \(R\text{-}\mathbf{Mod}\rightarrow S\text{-}\mathbf{Mod}\) 的函子。以 \(\varphi\) 替代 \(\varphi^{-1}\) 作反向构造，两次作用恢复原作用，两个函子严格互逆。

\ProseParagraphBreak
如果只要求模范畴等价，是否仍能恢复环同构？\chapref{b2:ch:16}将研究这一问题，并给出比环同构更宽的 \term[Morita-dengjia]{Morita 等价}[Morita equivalence]：它指 \(R\text{-}\mathbf{Mod}\) 与 \(S\text{-}\mathbf{Mod}\) 的范畴等价。那里还将证明矩阵环与原环的具体等价。

\ProseParagraphBreak
由\cref{b2:prop:equivalence-preserves-universal-objects}，模范畴等价保留零模、直和及直积所表达的映射性质，但正则模不是这个等价自动识别的指定对象。若等价把 \({}_RR\) 送到某个 \(S\) 模 \(P\)，它所对应的是 \(P\) 的自同态及其复合，并不保证 \(P\cong{}_SS\)。\cref{b2:prop:regular-hom-endomorphism}给出 \(\operatorname{End}_R({}_RR)\cong R^{\mathrm{op}}\)，其中反环来自右乘与复合次序；要在另一边用同一公式识别 \(S^{\mathrm{op}}\)，还须知道所取对象确实是正则模。因此，模范畴等价本身不允许直接把两个正则模的自同态环等同起来。\chapref{b2:ch:16}将研究哪些模能够代替正则模来描述整个模范畴。


\begin{chaptersummary}
范畴记录对象、态射和复合，函子必须同时处理对象与态射。Hom 对第一个变量反变、对第二个变量协变，这一方向约定决定了自然性方块中的箭头。自由模的映射对应不仅逐对象成立，还与原集合和目标模的映射相容。自然变换正是记录这种相容性的一族分量；分量逐个可逆才得到自然同构。双对偶的求值映射给出自然同构，而有限生成交换群的挠商虽逐个可以分裂，却不能自然地选择截面。

\ProseParagraphBreak
泛性质把构造表述为相容映射的唯一分解。模的拉回是直和中的相等条件子模，推出则是把两条公共来源的像识别起来的商模。比较两种构造时，应保留投影、嵌入或商映射，而不能只比较抽象同构类型；无限直积与坐标嵌入就会因延拓不唯一而失去余积的泛性质。范畴等价中，全性使目标态射能够提升，忠实性使提升后的等式能够核对，本质满使每个目标对象都能找到同构代表。因此等价保留零对象、积和余积等映射性质。模范畴等价却未指定正则模的对应方式，而环的乘法可以从正则模的自同态中恢复；这一区别引出 Morita 理论的问题。
\end{chaptersummary}

\BookExercises

\begin{exercise}\label{b2:ex:05-one-object-actions}
把群 \(G\) 看成只有一个对象的范畴 \(\mathcal G\)。证明函子 \(\mathcal G\rightarrow\mathbf{Set}\) 等价于一个集合上的左 \(G\) 作用；两个这样的函子之间的自然变换等价于等变映射。若目标换成 \(\mathbf{Vect}_k\)，相应的数据是什么？
\end{exercise}
\begin{solution}
唯一对象的像是集合 \(X\)，每个 \(g\in G\) 给出函数 \(F(g):X\rightarrow X\)。函子公理说明 \(F(1)=\operatorname{id}_X\)、\(F(gh)=F(g)F(h)\)，故 \(g\cdot x=F(g)(x)\) 满足左作用公理；\(F(g^{-1})\) 又是逆函数，所以每个 \(g\) 确实作用为置换。反向把左作用逐个读成这些函数，即得函子。

自然变换只有一个分量 \(a:X\rightarrow Y\)，条件为 \(aF(g)=H(g)a\)，逐元素写成 \(a(g\cdot x)=g\cdot a(x)\)，正是等变性。目标为向量空间时，各 \(F(g)\) 是可逆线性映射，故得到群同态 \(G\rightarrow\operatorname{GL}(V)\)，即线性表示；自然变换的分量是与全部群作用相容的线性映射。
\end{solution}

\begin{exercise}\label{b2:ex:05-posets-opposite}
设 \(P,Q\) 为偏序集。说明保序映射与其范畴之间的函子如何对应，并证明两个保序映射 \(F,G:P\rightarrow Q\) 之间存在自然变换 \(F\Rightarrow G\)，当且仅当每个 \(x\in P\) 都满足 \(F(x)\le G(x)\)。这种自然变换何时可逆？
\end{exercise}
\begin{solution}
态射 \(x\rightarrow y\) 存在当且仅当 \(x\le y\)，所以函子对态射的赋值存在恰好要求保序；一旦存在，它由唯一性完全确定，并自动保持复合与恒等。自然变换在 \(x\) 处需要态射 \(F(x)\rightarrow G(x)\)，即题述不等式。若所有这些态射存在，自然性图的两条路径有相同起终点，因而由偏序范畴中态射的唯一性相等。可逆要求两向不等式，等价于 \(F(x)=G(x)\) 对所有 \(x\) 成立，此时两个函子本身相等。
\end{solution}

\begin{exercise}\label{b2:ex:05-free-functor-map}
设 \(R\ne0\)，\(a:X\rightarrow Y\) 为函数，自由函子诱导 \(F(a):R^{(X)}\rightarrow R^{(Y)}\)。证明：\(F(a)\) 单射当且仅当 \(a\) 单射；\(F(a)\) 满射当且仅当 \(a\) 满射。写出 \(X=\{1,2,3\}\)、\(Y=\{u,v\}\)、\(a(1)=a(2)=u,a(3)=v\) 时的核与矩阵。
\end{exercise}
\begin{solution}
若 \(a\) 单射，不同基向量映到不同的标准基向量，因此像为零时每个原系数都为零。若 \(a(x)=a(x')\) 且 \(x\ne x'\)，非零向量 \(e_x-e_{x'}\) 在核中，其中非零性使用 \(1_R\ne0\)。若 \(a\) 满射，每个 \(e_y\) 都有原像，故 \(F(a)\) 满。反之，若 \(y\notin a(X)\)，所有像在 \(y\) 坐标上为零，不能得到 \(e_y\)。

具体矩阵是 \(\begin{pmatrix}1&1&0\\0&0&1\end{pmatrix}\)，核由 \((r,-r,0)\)、\(r\in R\) 组成，以 \(e_1-e_2\) 为基。矩阵条目均为 \(0,1\)，即使 \(R\) 非交换，所写映射也与左标量作用相容。
\end{solution}

\begin{exercise}\label{b2:ex:05-natural-center}
设 \(R\) 为含幺环，\(U:R\text{-}\mathbf{Mod}\rightarrow\mathbf{Ab}\) 为忘却到加法群的函子。证明 \(U\) 的自然自变换在逐分量相加和复合下组成与 \(R\) 同构的环，确定其中的自然自同构。哪些自然自变换的全部分量还是 \(R\) 模同态？用 \(R=M_2(k)\) 给出一个分量不是 \(R\) 线性的例子。
\end{exercise}
\begin{solution}
设 \(\eta:U\Rightarrow U\) 自然，令 \(c=\eta_R(1)\)。对任意 \(M\) 和 \(m\in M\)，映射 \(f_m:R\rightarrow M\)、\(r\mapsto rm\) 是 \(R\) 线性的，因而自然性在 \(1\) 处给出
\[
\eta_M(m)=U(f_m)\eta_R(1)=cm.
\]
这里 \(\eta_R\) 只须为加法群同态，不能预先使用它的 \(R\) 线性。反过来，任意 \(c\in R\) 所给的 \(m\mapsto cm\) 都保持加法，并由 \(f(cm)=cf(m)\) 与所有模同态相容。于是这些自然自变换恰与 \(R\) 的元素对应。逐分量相加和复合分别对应 \(c+d\) 和 \(cd\)，恒等变换对应 \(1\)；自然自同构恰对应 \(R^\times\)。

若全部分量还是 \(R\) 线性的，在正则模上便有 \(cr=\eta_R(r)=r\eta_R(1)=rc\)，所以 \(c\in Z(R)\)；中心元素也确实保证全部分量 \(R\) 线性，这与\cref{b2:prop:natural-center}一致。取 \(R=M_2(k)\)、\(c=E_{12}\)、\(r=E_{21}\)，则 \(cr=E_{11}\ne E_{22}=rc\)。因此左乘 \(E_{12}\) 给出忘却函子的自然自变换，却不是正则左模的自同态。
\end{solution}

\begin{exercise}\label{b2:ex:05-no-natural-torsion-splitting}
设 \(\mathcal G\) 的对象为有限生成交换群，态射只取它们之间的同构，\(J:\mathcal G\rightarrow\mathbf{Ab}_{\mathrm{fg}}\) 为包含函子。将挠商函子限制在这个定义域上，得到 \(Q:\mathcal G\rightarrow\mathbf{Ab}_{\mathrm{fg}}\)、\(Q(A)=A/T(A)\)，商映射给出自然变换 \(q:J\Rightarrow Q\)。证明不存在自然变换 \(s:Q\Rightarrow J\) 使每个 \(q_As_A=\operatorname{id}_{Q(A)}\)，即使只要求与同构相容，也不能自然地选择截面。具体地，对 \(m\ge2\) 和 \(A=\mathbb Z\oplus\mathbb Z/m\mathbb Z\)，求出 \(q_A\) 的全部截面，并考察自同构
\[
\alpha(n,\bar a)=(n,\bar a+\bar n)
\]
对这些截面的作用。
\end{exercise}
\begin{solution}
由\cref{b2:prop:no-natural-torsion-splitting}，只限制定义域的态射后，\(Q,q\) 仍满足各自的公理，各个商映射也仍有截面。这里的商映射和截面是目标范畴 \(\mathbf{Ab}_{\mathrm{fg}}\) 中的群同态，不要求属于 \(\mathcal G\)。在题设 \(A\) 上，挠子群为 \(0\oplus\mathbb Z/m\mathbb Z\)，故把 \(Q(A)\) 识别为 \(\mathbb Z\) 后，\(q_A(n,\bar a)=n\)。截面由 \(1\) 的像唯一确定；其第一坐标必须为 \(1\)，所以全部截面为
\[
s_\delta(n)=(n,n\delta),\qquad \delta\in\mathbb Z/m\mathbb Z.
\]
映射 \(\alpha\) 的逆为 \((n,\bar a)\mapsto(n,\bar a-\bar n)\)，且 \(Q(\alpha)=\operatorname{id}_{\mathbb Z}\)。但
\[
\alpha s_\delta=s_{\delta+\bar1},
\]
所以 \(\alpha\) 循环置换这 \(m\) 个截面，没有固定截面。若截面族自然，则在 \(\alpha\) 处必须满足 \(\alpha s_A=s_AQ(\alpha)=s_A\)，矛盾。这里的障碍仅由一个对象的自同构就已产生。
\end{solution}

\begin{exercise}\label{b2:ex:05-evaluation-dual}
对任意域上的向量空间 \(V\)，证明 \(j_V^*j_{V^*}=\operatorname{id}_{V^*}\)，其中 \(j\) 是双对偶求值映射。说明这个等式为何不要求 \(V\) 有限维，也不意味着 \(j_V\) 在无限维时满射。
\end{exercise}
\begin{solution}
对 \(\lambda\in V^*\)、\(v\in V\)，按两次求值的定义，
\[
\bigl(j_V^*j_{V^*}(\lambda)\bigr)(v)
=j_{V^*}(\lambda)\bigl(j_V(v)\bigr)
=j_V(v)(\lambda)=\lambda(v).
\]
因此等式成立，没有一步需要选基或有限维。它表明 \(j_{V^*}\) 有左逆，\(j_V^*\) 有右逆；它没有给出 \(j_V\) 的右逆。事实上\cref{b2:ex:infinite-double-dual}中 \(j_V\) 不满，而此等式仍成立。
\end{solution}

\begin{exercise}\label{b2:ex:05-initial-terminal-rings}
核对 \(\mathbf{Ring}\) 的初始对象与终对象。再考虑由所有域及保单位元同态组成的范畴，证明其中既没有初始对象，也没有终对象。
\end{exercise}
\begin{solution}
从 \(\mathbb Z\) 到环 \(R\) 的保单位元同态只能是 \(n\mapsto n1_R\)，且确实存在；从 \(R\) 到零环只有零映射，它保持单位元，因为零环的单位元等于零。因此分别得到初始对象和终对象。

域之间的保单位元同态单射：其核为理想，不能含 \(1\)，所以只能为零。因此这样的同态保持特征。若某个域初始，它必须同时有到 \(\mathbb F_2\) 和 \(\mathbb F_3\) 的同态，特征须同时为二和三，矛盾。若域 \(K\) 为终对象，则必须同时有 \(\mathbb F_2\rightarrow K\) 和 \(\mathbb F_3\rightarrow K\) 的同态；前者要求 \(\operatorname{char}K=2\)，后者要求 \(\operatorname{char}K=3\)，同样矛盾。
\end{solution}

\begin{exercise}\label{b2:ex:05-divisibility-products}
在正整数按整除关系组成的偏序范畴中，确定初始对象、两个对象的积与余积。是否有终对象？说明结论中的“积”为什么不是通常的整数乘积。
\end{exercise}
\begin{solution}
\(1\) 整除每个正整数，故为初始对象。积 \(p\) 要整除 \(a,b\)，并使任一共同因子 \(d\) 都整除 \(p\)，所以 \(p=\gcd(a,b)\)。余积 \(q\) 要被 \(a,b\) 整除，并整除任一公倍数，故 \(q=\operatorname{lcm}(a,b)\)。不存在被全部正整数整除的正整数，因而没有终对象。比如 \(2,4\) 的范畴积为二，通常乘积为八，后者甚至不能给出到二的态射。
\end{solution}

\begin{exercise}\label{b2:ex:05-infinite-product-not-coproduct}
设 \(k\) 为域，\(P=k^{\mathbb N}\)、\(S=k^{(\mathbb N)}\subseteq P\)，并以 \(e_n\) 表示标准坐标向量。证明限制映射
\[
\rho:P^*\longrightarrow k^{\mathbb N},\qquad
\ell\longmapsto\bigl(\ell(e_n)\bigr)_{n\in\mathbb N}
\]
满射而不单射，并构造 \(\ker\rho\cong(P/S)^*\) 的自然同构。解释为何指定所有标准坐标向量上的值，仍不能唯一确定 \(P\) 上的线性泛函。
\end{exercise}
\begin{solution}
给定 \((a_n)\in k^{\mathbb N}\)，公式 \(\ell_0(\sum_n x_ne_n)=\sum_n x_na_n\) 在 \(S\) 上定义线性泛函，因为每个向量仅有有限多个非零坐标。由\cref{b2:lem:functional-extension}，\(\ell_0\) 可延拓到 \(P\)，因此 \(\rho\) 满射。

\(\rho(\ell)=0\) 当且仅当 \(\ell\) 在所有 \(e_n\) 上为零，也就是 \(\ell|_S=0\)。记商映射为 \(\pi:P\rightarrow P/S\)，这样的 \(\ell\) 唯一写成 \(\lambda\pi\)，其中 \(\lambda:P/S\rightarrow k\) 为线性泛函：具体地，令 \(\lambda(v+S)=\ell(v)\)，由 \(\ell|_S=0\) 保证良定义。因此 \(\lambda\mapsto\lambda\pi\) 给出线性同构 \((P/S)^*\cong\ker\rho\)，且只使用商映射，没有选择。

向量 \(w=(1,1,\ldots)\) 不在 \(S\) 中，所以 \(w+S\ne0\)。先在 \(k(w+S)\) 上定义取值为 \(1\) 的泛函，再用泛函延拓引理延拓到 \(P/S\)，得到非零的 \(\lambda\)。于是 \(\lambda\pi\ne0\) 却属于 \(\ker\rho\)，故 \(\rho\) 不单射。标准坐标向量只张成 \(S\)，其上取值没有控制剩余的商空间；这正是\cref{b2:prop:infinite-product-not-coproduct}中唯一性失败的原因。
\end{solution}

\begin{exercise}\label{b2:ex:05-pullback-congruences}
把 \(\mathbb Z/6\mathbb Z\) 与 \(\mathbb Z/10\mathbb Z\) 都按自然商映射送到 \(\mathbb Z/2\mathbb Z\)。求其拉回，并给出拉回与 \(\mathbb Z/30\mathbb Z\) 的具体同构及两个投影。
\end{exercise}
\begin{solution}
由\cref{b2:prop:module-pullback-pushout}，拉回是两坐标具有相同奇偶性的剩余类对。对整数 \(t\)，记 \([t]_m\) 为它模 \(m\) 的剩余类。映射
\[
\psi:\mathbb Z/30\mathbb Z\longrightarrow
(\mathbb Z/6\mathbb Z)\times_{\mathbb Z/2\mathbb Z}(\mathbb Z/10\mathbb Z),
\quad [t]_{30}\longmapsto([t]_6,[t]_{10})
\]
良定义，因为三十同时被六和十整除。若像为零，则六和十都整除 \(t\)，故三十整除 \(t\)，所以 \(\psi\) 单射。拉回中，六个第一坐标各有五个同奇偶性的第二坐标，共三十个元素，因而单射已经满射。两个投影在此同构下就是模六和模十的商映射。
\end{solution}

\begin{exercise}\label{b2:ex:05-pushout-torsion}
求整数模同态 \(u,v:\mathbb Z\rightarrow\mathbb Z\)、\(u(t)=4t\)、\(v(t)=6t\) 的推出。除了给出抽象同构类型，还要写出两个结构映射在所选分解下的公式。
\end{exercise}
\begin{solution}
推出为 \(Q=\mathbb Z^2/\mathbb Z(4,-6)\)。向量 \(a=(2,-3)\)、\(b=(-1,2)\) 构成 \(\mathbb Z^2\) 的基，因为以它们为列的矩阵行列式为一。关系为 \(2a=0\)，故
\[
Q=\langle\bar a\rangle\oplus\langle\bar b\rangle
\cong\mathbb Z/2\mathbb Z\oplus\mathbb Z.
\]
由 \((1,0)=2a+3b\)、\((0,1)=a+2b\)，两个结构映射为
\[
i(t)=(\bar0,3t),\qquad j(t)=(\bar t,2t).
\]
确有 \(i(4t)=(\bar0,12t)=j(6t)\)。仅写 \(Q\cong\mathbb Z/2\mathbb Z\oplus\mathbb Z\) 会遗漏这两个映射，因此不足以单独记录原推出图。
\end{solution}

\begin{exercise}\label{b2:ex:05-pullback-surjection-pushout-injection}
在模的拉回图中，若 \(f:M\rightarrow P\) 满射，证明投影 \(q:M\times_PN\rightarrow N\) 满射，且其核自然同构于 \(\ker f\)。在模的推出图中，若 \(u:P\rightarrow M\) 单射，证明 \(j:N\rightarrow M\amalg_PN\) 单射。
\end{exercise}
\begin{solution}
给定 \(n\in N\)，由 \(f\) 满，取 \(m\in M\) 使 \(f(m)=g(n)\)。此时 \((m,n)\) 在拉回中并投影到 \(n\)。投影核由 \((m,0)\)、\(f(m)=0\) 组成，因此 \(m\mapsto(m,0)\) 给出所述同构。

若推出中 \(j(n)=0\)，式\eqref{b2:eq:module-pushout}说明 \((0,n)=\bigl(u(p),-v(p)\bigr)\) 对某个 \(p\) 成立。\(u\) 单射使 \(p=0\)，于是 \(n=0\)，所以 \(j\) 单射。此处必须使用关系子模正好等于同态 \(p\mapsto\bigl(u(p),-v(p)\bigr)\) 的像，而不是只知道其中含有这些关系。
\end{solution}

\begin{exercise}\label{b2:ex:05-regular-module-detects}
设 \(R\ne0\) 为含幺环，\(I\) 为集合，\(F=R^{(I)}\)。证明 \(\operatorname{Hom}_R(F,-)\) 能检测所有左 \(R\) 模同构，当且仅当 \(I\ne\varnothing\)。这里“检测”是指：对任意模同态 \(f:M\rightarrow N\)，若 \(\operatorname{Hom}_R(F,f)\) 为双射，则 \(f\) 为同构。
\end{exercise}
\begin{solution}
由自由模的泛性质，同态 \(F\rightarrow M\) 对应标准基向量的任意像族，所以有自然双射 \(\operatorname{Hom}_R(F,M)\cong M^I\)。在此识别下，后复合 \(f\) 对应逐坐标映射 \(f^I:M^I\rightarrow N^I\)。

若 \(I\ne\varnothing\)，固定 \(i_0\in I\)。当 \(f(m)=f(m')\) 时，把 \(m,m'\) 分别放在 \(i_0\) 坐标，其余坐标取零；\(f^I\) 的单射性使这两族相等，故 \(m=m'\)。给定 \(n\in N\)，把 \(n\) 放在 \(i_0\) 坐标、其余取零；\(f^I\) 的满射性给出一族原像，其 \(i_0\) 坐标就是 \(n\) 在 \(f\) 下的原像。因此 \(f\) 双射，且模同态的逆保持加法和标量乘法，故为模同构。反过来，同构经任何函子仍为同构，所以一定给出双射。这推广了\cref{b2:cor:regular-module-detects}。

若 \(I=\varnothing\)，则 \(F=0\)，从 \(F\) 到任意模的同态集合都是单元素集合。因此 \(0\rightarrow R\) 在 \(\operatorname{Hom}_R(F,-)\) 下给出双射，而它本身不是同构，因为 \(R\ne0\)。
\end{solution}

\begin{exercise}\label{b2:ex:05-equivalence-hypotheses}
分别给出如下函子：忠实且本质满但不全；全且本质满但不忠实；全忠实但不本质满。逐个指出它为何不是等价。
\end{exercise}
\begin{solution}
第一个取两元素反链到两元素链的恒等对象映射。它忠实且对象满，但目标中从较小对象到较大对象的态射无原像，所以不全。第二个取群 \(\mathbb Z/2\mathbb Z\) 所给的单对象范畴到平凡单对象范畴的函子。每个目标态射均有原像且唯一对象在像中，但两个不同群元素被合并，所以不忠实。第三个取只含零模的范畴到 \(\mathbb Z\text{-}\mathbf{Mod}\) 的包含函子。它在唯一 Hom 集合上为双射，却没有与 \(\mathbb Z\) 同构的像对象，故不本质满。\cref{b2:thm:equivalence-criterion}分别以所缺条件排除等价。
\end{solution}

\begin{exercise}\label{b2:ex:05-poset-equivalence}
证明两个偏序集作为范畴等价，当且仅当它们作为偏序集同构。再证明两个群作为单对象范畴等价，当且仅当这两个群同构。
\end{exercise}
\begin{solution}
对偏序集，全忠实函子给出 \(x\le y\iff F(x)\le F(y)\)。若 \(F(x)=F(y)\)，两向不等式推出 \(x=y\)，故对象作用单射。偏序范畴中的同构只可能是对象相等，所以本质满就是对象满射。因此 \(F\) 是双射并双向保持序，即偏序同构；反向显然产生范畴同构，故也是等价。

群的单对象范畴之间，函子的态射作用正是群同态，因为它保持乘法与单位元。唯一对象使本质满自动满足，全忠实则要求这个群同态双射。因此等价恰好给出群同构，反向同样成立。
\end{solution}

\begin{exercise}\label{b2:ex:05-coordinate-choice-natural}
为每个有限维 \(k\) 向量空间 \(V\) 选择三个坐标同构 \(c_V,d_V,e_V:k^{\dim V}\rightarrow V\)，得到矩阵函子 \(H_c,H_d,H_e\)。令 \(P_{dc}(V)=d_V^{-1}c_V\)，其余换基矩阵同样定义。证明
\[
P_{ed}(V)P_{dc}(V)=P_{ec}(V),\qquad
P_{cc}(V)=I,\qquad P_{cd}(V)=P_{dc}(V)^{-1}.
\]
在 \(V=k^2\) 上，取 \(c_V=I\)，并取 \(d_V,e_V\) 的矩阵为
\[
D=\begin{pmatrix}1&0\\1&1\end{pmatrix},\qquad
E=\begin{pmatrix}0&1\\1&0\end{pmatrix}.
\]
计算这三个换基矩阵，以及线性映射 \(f:V\rightarrow V\)、\(f(x,y)=(y,0)\) 在三组坐标下的矩阵，并核对换基的自然性等式。
\end{exercise}
\begin{solution}
由\cref{b2:prop:coordinate-choice-natural}，这些矩阵分别给出相应坐标函子间的自然同构。两个连续转换的分量满足
\[
(e_V^{-1}d_V)(d_V^{-1}c_V)=e_V^{-1}c_V,
\]
所以从 \(c\) 经 \(d\) 再到 \(e\)，与直接从 \(c\) 到 \(e\) 相同。其余两个等式由 \(c_V^{-1}c_V=I\) 和逆映射的公式得到。

在题设坐标下，
\[
P_{dc}=\begin{pmatrix}1&0\\-1&1\end{pmatrix},\qquad
P_{ed}=\begin{pmatrix}1&1\\1&0\end{pmatrix},\qquad
P_{ec}=\begin{pmatrix}0&1\\1&0\end{pmatrix},
\]
相乘确有 \(P_{ed}P_{dc}=P_{ec}\)。若 \(A=H_c(f)=\begin{pmatrix}0&1\\0&0\end{pmatrix}\)，则
\[
H_d(f)=D^{-1}AD=\begin{pmatrix}1&1\\-1&-1\end{pmatrix},
\qquad H_e(f)=E^{-1}AE=\begin{pmatrix}0&0\\1&0\end{pmatrix}.
\]
逐个相乘得到
\[
H_d(f)P_{dc}=P_{dc}A=\begin{pmatrix}0&1\\0&-1\end{pmatrix},
\qquad
H_e(f)P_{ed}=P_{ed}H_d(f)=\begin{pmatrix}0&0\\1&1\end{pmatrix}.
\]
这核对了两次换基的自然性；合成后亦有 \(H_e(f)P_{ec}=P_{ec}A\)。等式在任意特征下都成立，其中负号按域 \(k\) 的运算解释。
\end{solution}

\begin{exercise}\label{b2:ex:05-equivalence-kernels}
设 \(F:R\text{-}\mathbf{Mod}\rightarrow S\text{-}\mathbf{Mod}\) 是范畴等价。若 \(k:K\rightarrow M\) 为同态 \(f:M\rightarrow N\) 的核嵌入，证明 \(F(k)\) 具有 \(F(f)\) 的核的泛性质。证明时不要假设 \(F\) 保持底集合。
\end{exercise}
\begin{solution}
由\cref{b2:prop:equivalence-preserves-universal-objects}，\(F\) 把零对象送到零对象；任何零态射都经过零对象分解，故其像仍是零态射。因此 \(F(f)F(k)=0\)。给定 \(a:Y\rightarrow F(M)\) 且 \(F(f)a=0\)，由本质满选同构 \(e:F(T)\rightarrow Y\)，再由全忠实性唯一提升 \(ae\) 为 \(b:T\rightarrow M\)。忠实性将 \(F(fb)=0\) 送回 \(fb=0\)，因此存在唯一 \(c:T\rightarrow K\) 使 \(kc=b\)。

令 \(h=F(c)e^{-1}\)，则 \(F(k)h=a\)。若另有 \(h'\) 满足，把 \(h'e\) 唯一提升为 \(c':T\rightarrow K\)。忠实性给出 \(kc'=b\)，由核的唯一性得 \(c'=c\)，所以 \(h'=h\)。这里由 \(fb=0\) 得到 \(\operatorname{im}b\subseteq\ker f\)，再由核嵌入 \(k:K\rightarrow M\) 的泛性质得到唯一分解；不需要假设 \(F\) 保持底集合。
\end{solution}

\begin{exercise}\label{b2:ex:05-relabel-scalars}
设 \(\varphi:R\rightarrow S\) 为环同构，按正文把左 \(R\) 模变成左 \(S\) 模。证明这个函子全忠实，并直接写出逆函子。说明只假设 \(\varphi\) 为环同态时，为什么原公式不能照用，以及哪一个方向仍有无条件的忘却标量函子。
\end{exercise}
\begin{solution}
对给定 \(M,N\)，加法群同态 \(f:M\rightarrow N\) 保持新 \(S\) 作用，当且仅当对所有 \(s\in S\) 有 \(f\bigl(\varphi^{-1}(s)m\bigr)=\varphi^{-1}(s)f(m)\)。\(\varphi^{-1}\) 满到 \(R\)，这恰等价于 \(R\) 线性。因此每个 Hom 集合原样对应，函子全忠实。逆函子将左 \(S\) 模 \(N\) 的作用改写成 \(r\cdot n=\varphi(r)n\)，两次改写严格恢复原作用。

一般环同态没有可用的 \(\varphi^{-1}\)，所以不能由原公式把每个 \(R\) 模变成 \(S\) 模。但对任意左 \(S\) 模，总可用 \(r\cdot n=\varphi(r)n\) 得到左 \(R\) 模；原 \(S\) 线性映射仍 \(R\) 线性，因而得到从 \(S\text{-}\mathbf{Mod}\) 到 \(R\text{-}\mathbf{Mod}\) 的函子。反方向的标量扩张需要第六章的张量积。
\end{solution}
